% ============================================================
%  HW02 – Domain Testing on EShop | Report Template
% ============================================================
\documentclass[12pt, a4paper]{report}

% ---------- Encoding & Language ----------
\usepackage[utf8]{inputenc}
\usepackage[T5]{fontenc}
\usepackage[vietnamese, english]{babel}

% ---------- Page Layout ----------
\usepackage[
    top=2.5cm, bottom=2.5cm,
    left=3cm,  right=2.5cm
]{geometry}

% ---------- Typography ----------
\usepackage{lmodern}
\usepackage{microtype}
\usepackage{setspace}
\onehalfspacing

% ---------- Colors (đen trắng / grayscale only) ----------
\usepackage[dvipsnames, table]{xcolor}
% Toàn bộ báo cáo dùng thuần đen trắng: mọi màu sắc được ép về đen/trắng/xám.
\definecolor{primaryblue}{gray}{0}          % tiêu đề, link, caption -> đen
\definecolor{headerblue}{gray}{0}
\definecolor{lightgray}{gray}{0.95}         % nền code -> xám rất nhạt
\definecolor{pastelgreen}{gray}{1}          % nền box -> trắng
\definecolor{pastelred}{gray}{1}
\definecolor{pastelyellow}{gray}{1}
% Ép các màu chuẩn (dùng trong \textcolor Pass/Fail/Note, code) về đen:
\definecolor{green}{gray}{0}
\definecolor{red}{gray}{0}
\definecolor{orange}{gray}{0}
\definecolor{BrickRed}{gray}{0}

% ---------- Math ----------
\usepackage{amsmath, amssymb}

% ---------- Tables ----------
\usepackage{booktabs}
\usepackage{longtable}
\usepackage{tabularx}
\usepackage{multirow}
\usepackage{array}
\newcolumntype{L}[1]{>{\raggedright\arraybackslash}p{#1}}
\newcolumntype{C}[1]{>{\centering\arraybackslash}p{#1}}
\newcolumntype{R}[1]{>{\raggedleft\arraybackslash}p{#1}}

% ---------- Graphics ----------
\usepackage{graphicx}
\usepackage{float}
\graphicspath{{./figures/}}

% ---------- Listings (Code) ----------
\usepackage{listings}
\lstset{
    basicstyle=\ttfamily\small,
    keywordstyle=\color{primaryblue}\bfseries,
    commentstyle=\color{gray}\itshape,
    stringstyle=\color{BrickRed},
    backgroundcolor=\color{lightgray},
    frame=single,
    rulecolor=\color{gray!40},
    breaklines=true,
    showstringspaces=false,
    numbers=left,
    numberstyle=\tiny\color{gray},
    tabsize=4,
    captionpos=b
}

% ---------- Hyperlinks ----------
\usepackage[
    colorlinks=true,
    linkcolor=primaryblue,
    urlcolor=primaryblue,
    citecolor=primaryblue,
    pdftitle={HW02 – Domain Testing on EShop},
    pdfauthor={Lý Quốc Thạnh},
    pdfsubject={Software Testing}
]{hyperref}

% ---------- Headers & Footers ----------
\usepackage{fancyhdr}
\pagestyle{fancy}
\fancyhf{}
\fancyhead[L]{\small\textcolor{gray}{\subjectName}}
\fancyhead[R]{\small\textcolor{gray}{\hwTitle}}
\fancyfoot[C]{\small\thepage}
\renewcommand{\headrulewidth}{0.4pt}
\renewcommand{\footrulewidth}{0pt}

% ---------- Section Styling ----------
\usepackage{titlesec}
\titleformat{\chapter}[block]
    {\normalfont\LARGE\bfseries\color{primaryblue}}
    {\thechapter.}{1em}{}
\titleformat{\section}
    {\normalfont\Large\bfseries\color{primaryblue!80}}
    {\thesection}{1em}{}
\titleformat{\subsection}
    {\normalfont\large\bfseries}
    {\thesubsection}{1em}{}
\titleformat{\subsubsection}
    {\normalfont\normalsize\bfseries}
    {\thesubsubsection}{1em}{}

% ---------- Captions ----------
\usepackage{caption}
\captionsetup{
    font=small,
    labelfont={bf, color=primaryblue},
    skip=6pt
}

% ---------- Enumitem ----------
\usepackage{enumitem}
\setlist[itemize]{leftmargin=*, topsep=4pt, itemsep=2pt}
\setlist[enumerate]{leftmargin=*, topsep=4pt, itemsep=2pt}

% ---------- Boxes & Notes ----------
\usepackage{tcolorbox}
\tcbuselibrary{skins, breakable}
\newtcolorbox{notebox}[1][]{
    colback=white, colframe=black, coltitle=black, colbacktitle=lightgray,
    title={\textbf{Note}}, fonttitle=\bfseries,
    #1
}
\newtcolorbox{bugbox}[1][]{
    colback=white, colframe=black, coltitle=black, colbacktitle=lightgray,
    title={\textbf{Bug}}, fonttitle=\bfseries,
    #1
}
\newtcolorbox{passbox}[1][]{
    colback=white, colframe=black, coltitle=black, colbacktitle=lightgray,
    title={\textbf{Pass}}, fonttitle=\bfseries,
    #1
}

% ---------- Metadata ----------
\newcommand{\subjectName}{KIỂM CHỨNG PHẦN MỀM}
\newcommand{\hwTitle}{HW02 – Domain Testing on EShop}
\newcommand{\instructor}{%
    \begin{tabular}[t]{@{}l@{}}
        Lâm Quang Vũ \\
        Trương Phước Lộc \\
        Hồ Tuấn Thanh
    \end{tabular}%
}
\newcommand{\studentName}{Lý Quốc Thạnh}
\newcommand{\studentID}{23127262}
\newcommand{\studentEmail}{lqthanh23@clc.fitus.edu.vn}

% ============================================================
\begin{document}

% ===== TRANG BÌA =====
\begin{titlepage}
    \centering
    \vspace*{1cm}
    \textbf{\Large TRƯỜNG ĐẠI HỌC KHOA HỌC TỰ NHIÊN}\\
    \textbf{\large KHOA CÔNG NGHỆ THÔNG TIN}\\
    \vspace{3cm}
    {\LARGE \textbf{\subjectName}}\\[2cm]
    {\LARGE \textbf{\hwTitle}}\\[2cm]
    \vspace{1.5cm}
    \begin{table}[h]
        \centering
        \large
        \begin{tabular}{ll}
            \textbf{Giảng viên:}  & \instructor \\[0.5cm]
            \textbf{Sinh viên:}   & \studentName \\
            \textbf{MSSV:}        & \studentID \\
            \textbf{Email:}       & \studentEmail \\
        \end{tabular}
    \end{table}
    \vfill
    {\large TP. Hồ Chí Minh, \today}
\end{titlepage}

% ===== MỤC LỤC =====
\tableofcontents
\listoffigures
\listoftables
\newpage

% ============================================================
%  CHƯƠNG 1: TỔNG QUAN
% ============================================================
\chapter{Tổng Quan}

\section{Các Tính Năng Được Chọn}
\begin{table}[H]
    \centering
    \caption{Tính năng được chọn}
    \begin{tabularx}{\linewidth}{C{1.5cm} C{2cm} L{4cm} L{6cm}}
        \toprule
        \textbf{Pool} & \textbf{FR-ID} & \textbf{Tên tính năng} & \textbf{Mô tả} \\
        \midrule
        A & FR-01 & Đăng ký tài khoản & Tạo tài khoản mới với họ tên, email, mật khẩu. \\
        B & FR-09 & Mã giảm giá & Áp dụng mã giảm giá (percent/fixed) cho đơn hàng. \\
        C & FR-14 & Quản lý danh mục (CRUD) & Thêm/xóa danh mục sản phẩm trên ứng dụng admin (\texttt{frontend-admin}, tab ``Danh mục''); UI chỉ có Thêm/Xóa, không có Sửa. \\
        D & FR-07 & Giỏ hàng (Mobile) & Thêm/sửa số lượng/xóa sản phẩm trong giỏ trên app mobile. \\
        \bottomrule
    \end{tabularx}
\end{table}

% ============================================================
%  CHƯƠNG 2: FEATURE A — POOL A
% ============================================================
\chapter{Feature A — Đăng ký tài khoản (Pool A)}

\section{Mô Tả Tính Năng}
Tính năng FR-01 — Đăng ký tài khoản cho phép người dùng tạo tài khoản mới bằng
họ tên, email và mật khẩu. Nó được cài đặt qua hai tầng:

\begin{itemize}
    \item \textbf{Tầng giao diện (UI):} \texttt{frontend-web/src/pages/Register.jsx} — biểu mẫu
    với ba trường. Đây là nơi duy nhất có kiểm tra dữ liệu: một biểu thức chính quy
    (regex) kiểm tra ``mật khẩu mạnh'' phía client. Khi hợp lệ, nó gửi \texttt{POST} tới API
    rồi chuyển sang trang \texttt{/login}.
    \item \textbf{Tầng API:} \texttt{POST /api/register} (\texttt{backend/server.js}) — chèn
    thẳng \verb|{name, email, password}| vào bảng \texttt{users}, không kiểm tra hợp lệ,
    không kiểm tra trùng lặp, và lưu mật khẩu ở dạng văn bản thuần (plaintext).
\end{itemize}

\noindent\textbf{Đầu vào (input):} ba ô của biểu mẫu Đăng ký — \texttt{name}, \texttt{email},
\texttt{password}.\\
\textbf{Đầu ra (output) quan sát trên UI:} thông báo lỗi ngay dưới form (vd ``Mật khẩu quá
yếu!''), hoặc đăng ký thành công và điều hướng sang \texttt{/login}. (Việc email trùng
có thể đối chiếu thêm ở tab ``Người dùng'' của app admin.)

\vspace{4pt}
\noindent\textbf{Ràng buộc mật khẩu theo đặc tả} (chữ trợ giúp trong \texttt{Register.jsx}):
``Tối thiểu 8 ký tự, có chữ hoa, chữ thường, số và ký tự đặc biệt.''
Tuy nhiên regex thực tế được cài đặt là:

\begin{lstlisting}[caption={Regex kiểm tra mật khẩu (Register.jsx)},numbers=none]
/^(?=.*[a-z])(?=.*[A-Z])(?=.*\d)(?=.*\s)[A-Za-z\d\s]{8,}$/
\end{lstlisting}

\begin{notebox}
Regex này bắt buộc có khoảng trắng (\texttt{(?=.*\textbackslash s)}) và, do lớp ký tự
là \texttt{[A-Za-z\textbackslash d\textbackslash s]}, nó cấm mọi ký tự đặc biệt — ngược
hoàn toàn với đặc tả (vốn yêu cầu phải có ký tự đặc biệt). Đây là khiếm khuyết trọng tâm
mà chương này nhắm tới. Hệ quả: mật khẩu mạnh như \texttt{Password123!} bị từ chối, còn
mật khẩu yếu như \texttt{Passw0rd 1} lại được chấp nhận.
\end{notebox}

\begin{notebox}
\textbf{Phạm vi \& quy ước cột kết quả.} Theo giảng viên (forum Q\&A HW02): functional
testing thực hiện từ UI Frontend. Mọi ca dưới đây thao tác trên biểu mẫu Đăng ký của
web (\texttt{frontend-web}, trang \texttt{/register}); cột Kết quả/Trạng thái là
hành vi quan sát trên giao diện (thông báo lỗi, điều hướng \texttt{/login}, và đối
chiếu bản ghi qua tab Người dùng của app admin). Chế độ ``sinh viên tự chạy, AI dự đoán từ mã
nguồn''. Cột Expected là hành vi đúng theo đặc tả; ca có Kết quả khác Expected là
lỗi.
\end{notebox}

\section{Domain Testing}

Domain Testing được áp dụng theo trình tự: (1) xác định biến $\to$
(2) phân hoạch miền thành lớp tương đương $\to$ (3) chọn điểm kiểm thử ON/OFF/IN/OUT trên biên
$\to$ (4) thiết kế ca. Nguyên tắc cốt lõi: coi mỗi biến đầu vào là một miền giá trị,
chia miền thành các lớp cho cùng một hành vi, rồi chỉ cần một điểm đại diện cho mỗi lớp
cộng với các điểm sát biên --- nhờ đó giảm số ca mà vẫn phủ hết hành vi. Với FR-01, biến trọng
tâm là \texttt{password} (nhiều lớp ký tự) và \texttt{email} (định dạng + tính duy nhất); biến
có thứ tự duy nhất để làm biên số học là \texttt{password.length}. Bước~5 tổng hợp ma trận truy
vết để chứng minh mọi lớp tương đương đều có ít nhất một ca kiểm thử.

\subsection{Bước 1: Xác Định Các Biến Đầu Vào}
Liệt kê mọi biến đầu vào (kể cả biến ẩn: độ dài mật khẩu, các lớp ký tự, tính duy nhất của
email).

\begin{table}[H]
    \centering
    \caption{Danh sách biến đầu vào — Feature A (FR-01)}
    \begin{tabularx}{\linewidth}{C{0.5cm} L{3cm} L{2.4cm} L{3.2cm} X}
        \toprule
        \textbf{\#} & \textbf{Biến} & \textbf{Kiểu dữ liệu} & \textbf{Ràng buộc} & \textbf{Ghi chú} \\
        \midrule
        1 & name & string (TEXT) & UI \texttt{required} & Không giới hạn độ dài/định dạng; backend không kiểm. \\
        2 & email & string (TEXT) & UI \texttt{required} & \texttt{type="text"} nên không kiểm định dạng; cột không \texttt{UNIQUE}. \\
        3 & password & string (TEXT) & regex client & Backend không kiểm; lưu plaintext. \\
        4 & password.length & integer (ẩn) & $\ge 8$ (regex \texttt{\{8,\}}) & Biến có thứ tự $\Rightarrow$ dùng cho BVA. \\
        5 & has\_lower / has\_upper / has\_digit & boolean (ẩn) & phải có mỗi loại & Điều kiện lớp ký tự của regex. \\
        6 & has\_whitespace & boolean (ẩn) & regex bắt buộc & Khiếm khuyết (đặc tả không yêu cầu). \\
        7 & has\_special\_char & boolean (ẩn) & regex cấm & Khiếm khuyết (đặc tả yêu cầu). \\
        8 & email.isUnique & boolean (ẩn) & phải duy nhất & Không nơi nào kiểm tra. \\
        \bottomrule
    \end{tabularx}
\end{table}

\subsection{Bước 2: Phân Hoạch Miền (Domain Partitioning)}
Mỗi biến được chia thành các lớp tương đương hợp lệ/không hợp lệ (có mã EC). Dấu *
đánh dấu lớp mà một hệ thống đúng phải xử lý theo đặc tả nhưng mã hiện tại làm sai.

\subsubsection{Biến name}
\begin{table}[H]
    \centering
    \caption{Phân hoạch miền — name}
    \begin{tabularx}{\linewidth}{C{1.4cm} L{6cm} C{1.6cm} X}
        \toprule
        \textbf{EC} & \textbf{Điều kiện} & \textbf{Loại} & \textbf{Ví dụ} \\
        \midrule
        EC-N1 & Tên thường, không rỗng & Valid & \texttt{Nguyen Van A} \\
        EC-N2 & Rỗng & Invalid & \verb|""| \\
        EC-N3 & Chỉ khoảng trắng & Invalid & \verb|"   "| \\
        EC-N4 & Rất dài (không có max) & Valid* & 100.000 ký tự \\
        \bottomrule
    \end{tabularx}
\end{table}

\subsubsection{Biến email}
\begin{table}[H]
    \centering
    \caption{Phân hoạch miền — email}
    \begin{tabularx}{\linewidth}{C{1.4cm} L{6cm} C{1.6cm} X}
        \toprule
        \textbf{EC} & \textbf{Điều kiện} & \textbf{Loại} & \textbf{Ví dụ} \\
        \midrule
        EC-E1 & Đúng định dạng, chưa dùng & Valid & \texttt{new@domain.com} \\
        EC-E2 & Sai định dạng & Invalid & \texttt{notanemail} \\
        EC-E3 & Rỗng & Invalid & \verb|""| \\
        EC-E4 & Trùng email đã tồn tại & Invalid & \texttt{test@eshop.com} \\
        EC-E5 & Có khoảng trắng bao quanh & Invalid & \verb|" test@eshop.com "| \\
        EC-E6 & Payload dạng SQL-injection & Invalid & \texttt{x'; DROP TABLE users;-{}-} \\
        \bottomrule
    \end{tabularx}
\end{table}

\subsubsection{Biến password}
\begin{table}[H]
    \centering
    \caption{Phân hoạch miền — password}
    \begin{tabularx}{\linewidth}{C{1.4cm} L{6cm} C{1.6cm} X}
        \toprule
        \textbf{EC} & \textbf{Điều kiện} & \textbf{Loại} & \textbf{Ví dụ} \\
        \midrule
        EC-P1 & Qua được regex (có ws, không ký tự đặc biệt, $\ge 8$) & Valid* & \texttt{Passw0rd 1} \\
        EC-P2 & Mạnh đúng đặc tả (hoa+thường+số+đặc biệt, không ws) & Spec-valid & \texttt{Password123!} \\
        EC-P3 & Thiếu chữ thường & Invalid & \texttt{PASSW0RD 1} \\
        EC-P4 & Thiếu chữ hoa & Invalid & \texttt{passw0rd 1} \\
        EC-P5 & Thiếu chữ số & Invalid & \texttt{Password x} \\
        EC-P6 & Thiếu khoảng trắng (điều kiện lỗi) & Invalid & \texttt{Password1} \\
        EC-P7 & Có ký tự đặc biệt & Invalid & \texttt{Passw0rd !} \\
        EC-P8 & Độ dài $<8$ & Invalid & \texttt{Aa1 bcd} \\
        EC-P9 & Rỗng & Invalid & \verb|""| \\
        \bottomrule
    \end{tabularx}
\end{table}

\noindentSpec-valid = đúng đặc tả nhưng bị mã từ chối; Valid* = mã
chấp nhận dù sai đặc tả. Hai lớp EC-P1 và EC-P2 chính là nơi lỗi biểu hiện rõ nhất.

\subsection{Bước 3: Chọn Điểm Kiểm Thử (Test Point Selection)}
Biến có thứ tự duy nhất là \texttt{password.length} với ranh giới $\ge 8$. Áp dụng chiến lược
ON/OFF/IN, mọi giá trị dưới đây giữ nguyên các điều kiện lớp ký tự còn lại (1 chữ thường, 1
chữ hoa, 1 chữ số, 1 khoảng trắng) để độ dài là biến duy nhất thay đổi.

\begin{table}[H]
    \centering
    \caption{Điểm ON/OFF/IN/OUT — password.length ($\ge 8$)}
    \begin{tabularx}{\linewidth}{L{2.4cm} L{3.2cm} C{2cm} X}
        \toprule
        \textbf{Loại điểm} & \textbf{Giá trị} & \textbf{Độ dài} & \textbf{Kỳ vọng} \\
        \midrule
        ON-point  & \texttt{Aa1 bcde}     & 8  & Chấp nhận (đúng biên) \\
        OFF-point & \texttt{Aa1 bcd}      & 7  & Từ chối (dưới biên) \\
        IN-point  & \texttt{Aa1 bcdefgh}  & 11 & Chấp nhận (trong miền) \\
        OUT-point & \verb|""|             & 0  & Từ chối (ngoài miền) \\
        \bottomrule
    \end{tabularx}
\end{table}

\noindent\textbf{Nhận xét quan trọng:} biên độ dài (\texttt{\{8,\}}) được cài đặt đúng
— không có off-by-one tại 8. Các khiếm khuyết nằm ở điều kiện lớp ký tự (bắt buộc
khoảng trắng, cấm ký tự đặc biệt), là các phân hoạch (Bước~2) chứ không phải biên số
học. Vì vậy Domain Testing phát hiện nhiều lỗi hơn BVA thuần túy ở tính năng này.

\subsection{Bước 4: Thiết Kế Test Cases — Domain Testing}
\begin{longtable}{C{1cm} L{3cm} L{3.5cm} L{3cm} C{2cm} C{1.5cm}}
    \caption{Test cases — Domain Testing, Feature A} \\
    \toprule
    \textbf{TC\#} & \textbf{Mô tả} & \textbf{Input} & \textbf{Expected Output} & \textbf{Kết quả} & \textbf{Trạng thái} \\
    \midrule
    \endfirsthead
    \multicolumn{6}{c}{\small\itshape (tiếp theo)} \\
    \toprule
    \textbf{TC\#} & \textbf{Mô tả} & \textbf{Input} & \textbf{Expected Output} & \textbf{Kết quả} & \textbf{Trạng thái} \\
    \midrule
    \endhead
    \midrule
    \multicolumn{6}{r}{\small\itshape (còn tiếp)} \\
    \endfoot
    \bottomrule
    \endlastfoot
    TC-A-DT-01 & MK mạnh có ký tự đặc biệt, đúng đặc tả \newline (BUG-A-01) & Form: name/email hợp lệ, pw=\texttt{Password123!}, bấm Đăng Ký & Đăng ký OK, sang \texttt{/login} & Báo ``MK quá yếu'' & \textcolor{red}{Fail} \\
    TC-A-DT-02 & MK qua regex: có ws, không ký tự đặc biệt \newline (BUG-A-02) & Form: pw=\texttt{Passw0rd 1} & Từ chối (thiếu ký tự đặc biệt) & Đăng ký OK, sang \texttt{/login} & \textcolor{red}{Fail} \\
    TC-A-DT-03 & Thiếu chữ hoa & Form: pw=\texttt{passw0rd 1} & Từ chối & Báo ``MK quá yếu'' & \textcolor{green!60!black}{Pass} \\
    TC-A-DT-04 & Thiếu chữ số & Form: pw=\texttt{Password x} & Từ chối & Báo ``MK quá yếu'' & \textcolor{green!60!black}{Pass} \\
    TC-A-DT-05 & Có ký tự đặc biệt kèm ws \newline (BUG-A-01) & Form: pw=\texttt{Passw0rd !} & Chấp nhận (MK mạnh) & Từ chối (cấm `!') & \textcolor{red}{Fail} \\
    TC-A-DT-06 & Email sai định dạng \newline (BUG-A-03) & Form: email=\texttt{notanemail}, pw hợp lệ & Từ chối email không hợp lệ & Đăng ký OK, sang \texttt{/login} & \textcolor{red}{Fail} \\
    TC-A-DT-07 & Email trùng đã tồn tại \newline (BUG-A-04) & Form: email=\texttt{test@eshop.com} (đã có) & Từ chối trùng & Đăng ký OK; tab admin có 2 dòng cùng email & \textcolor{red}{Fail} \\
    TC-A-DT-08 & Ký tự đặc biệt/injection ở email & Form: email=\texttt{x'; DROP TABLE users;-{}-} & Xử lý an toàn & Lưu literal, DB an toàn (tham số hóa) & \textcolor{green!60!black}{Pass} \\
    TC-A-DT-09 & Tên chỉ khoảng trắng & Form: name=\verb|"   "| & Từ chối (rỗng ngữ nghĩa) & \texttt{required} cho qua, đăng ký OK & \textcolor{red}{Fail} \\
\end{longtable}

\subsection{Bước 5: Ma Trận Truy Vết Độ Phủ (EC $\to$ Test Case)}
Bảng dưới đối chiếu mọi lớp tương đương ở Bước~2 với ca kiểm thử phủ nó (kể cả ca BVA),
chứng minh tiêu chí ``mỗi lớp $\ge 1$ điểm đại diện''. Một số lớp cùng họ hành vi dùng
chung một ca đại diện theo đúng nguyên tắc phân hoạch tương đương.

\begin{table}[H]
    \centering
    \caption{Truy vết EC $\to$ TC — Feature A (FR-01)}
    \begin{tabularx}{\linewidth}{C{1.8cm} C{3.4cm} X}
        \toprule
        \textbf{EC} & \textbf{Ca phủ (DT/BVA)} & \textbf{Ghi chú} \\
        \midrule
        EC-N1 & (nền ca hợp lệ) & Tên hợp lệ dùng ở TC-A-DT-01/02. \\
        EC-N2 & (UI \texttt{required} chặn) & Ô name rỗng không gửi được từ form --- hành vi đúng. \\
        EC-N3 & TC-A-DT-09 & Tên chỉ khoảng trắng. \\
        EC-N4 & TC-A-BV-05 & Tên rất dài (biên trên khuyết). \\
        EC-E1 & (nền ca hợp lệ) & Email hợp lệ dùng ở TC-A-DT-01/02. \\
        EC-E2 & TC-A-DT-06 & Sai định dạng. \\
        EC-E3 & (UI \texttt{required} chặn) & Ô email rỗng không gửi được từ form --- hành vi đúng. \\
        EC-E4 & TC-A-DT-07 & Email trùng. \\
        EC-E5 & TC-A-DT-07 (đại diện) & Biến thể khoảng trắng của lỗi định danh trùng; chi tiết ở Gap Analysis. \\
        EC-E6 & TC-A-DT-08 & Ký tự đặc biệt/injection. \\
        EC-P1 & TC-A-DT-02, BV-02 & Qua regex (ws, không đặc biệt, $\ge 8$). \\
        EC-P2 & TC-A-DT-01 & Mạnh đúng đặc tả. \\
        EC-P3 & TC-A-DT-03/04 (đại diện) & Cùng họ ``thiếu một lớp ký tự bắt buộc''; regex đối xứng nên một đại diện đủ. \\
        EC-P4 & TC-A-DT-03 & Thiếu chữ hoa. \\
        EC-P5 & TC-A-DT-04 & Thiếu chữ số. \\
        EC-P6 & TC-A-DT-01 & \texttt{Password123!} không có khoảng trắng $\Rightarrow$ bị từ chối (chính điều kiện EC-P6). \\
        EC-P7 & TC-A-DT-05 & Có ký tự đặc biệt. \\
        EC-P8 & TC-A-BV-01 & Độ dài $<8$. \\
        EC-P9 & TC-A-BV-04 & Rỗng (UI \texttt{required}+regex chặn). \\
        \bottomrule
    \end{tabularx}
\end{table}
\noindent\textbf{Kết luận độ phủ:} mọi lớp \texttt{name} (4), \texttt{email} (6), \texttt{password}
(9) đều có ca phủ, hoặc được UI \texttt{required} chặn ngay (lớp rỗng của name/email --- chính là
hành vi đúng, không gửi được qua form).

\section{Boundary Value Analysis (BVA)}

Boundary Value Analysis bổ trợ Domain Testing bằng cách tập trung vào vùng sát ranh giới
--- nơi lỗi hay xuất hiện nhất. Quy trình gồm: (1) xác định mọi biên (số học /
độ dài); (2) với mỗi biên sinh bộ ba giá trị \texttt{off (biên$-1$), on (biên), on$+1$}, cộng
điểm ngoài miền khi cần. Với FR-01, biên số học duy nhất là \texttt{password.length}$\,\ge 8$;
\texttt{name}/\texttt{email} chỉ có biên ``khác rỗng'' (một điểm: rỗng vs không rỗng) và
không có biên trên --- bản thân việc thiếu biên trên là một phát hiện. Bước~3 tổng hợp ma
trận phủ biên để chứng minh mỗi biên đều đủ điểm off/on/on$+1$.

\subsection{Bước 1: Xác Định Ranh Giới (Boundaries)}
Chỉ có \texttt{password.length} là biên số học được xác định rõ ($\ge 8$). \texttt{name} và
\texttt{email} chỉ có biên dưới ``khác rỗng'' (do UI \texttt{required}); không có biên
trên nào được đặc tả hay cưỡng chế (SQLite \texttt{TEXT} không giới hạn) — bản thân điều này
là một phát hiện (xem Gap Analysis).

\begin{table}[H]
    \centering
    \caption{Bảng giá trị biên — Feature A (FR-01)}
    \begin{tabularx}{\linewidth}{L{3cm} C{1.6cm} C{1.6cm} C{1.8cm} C{2.4cm} X}
        \toprule
        \textbf{Biến} & \textbf{Min} & \textbf{Min+1} & \textbf{Max-1} & \textbf{Max} & \textbf{Ghi chú} \\
        \midrule
        password.length & 8 & 9 & --- & --- (không có max) & Biên dưới đúng; off-point $=7$; không có biên trên. \\
        name.length     & 1 & 2 & --- & --- (không có max) & Chỉ \texttt{required} chặn rỗng ở UI. \\
        email.length    & 1 & 2 & --- & --- (không có max) & Tương tự; backend nhận cả \texttt{NULL}. \\
        \bottomrule
    \end{tabularx}
\end{table}

\subsection{Bước 2: Thiết Kế Test Cases — BVA}
\begin{longtable}{C{1cm} L{3cm} L{3.5cm} L{3cm} C{2cm} C{1.5cm}}
    \caption{Test cases — BVA, Feature A} \\
    \toprule
    \textbf{TC\#} & \textbf{Mô tả} & \textbf{Input} & \textbf{Expected Output} & \textbf{Kết quả} & \textbf{Trạng thái} \\
    \midrule
    \endfirsthead
    \toprule
    \textbf{TC\#} & \textbf{Mô tả} & \textbf{Input} & \textbf{Expected Output} & \textbf{Kết quả} & \textbf{Trạng thái} \\
    \midrule
    \endhead
    \bottomrule
    \endlastfoot
    TC-A-BV-01 & Min$-1$: độ dài 7 & Form: pw=\texttt{Aa1 bcd} (7) & Từ chối (dưới biên) & Báo ``MK quá yếu'' & \textcolor{green!60!black}{Pass} \\
    TC-A-BV-02 & Min: độ dài 8 & Form: pw=\texttt{Aa1 bcde} (8) & Chấp nhận (qua regex) & Đăng ký OK & \textcolor{green!60!black}{Pass} \\
    TC-A-BV-03 & Min$+1$: độ dài 9 & Form: pw=\texttt{Aa1 bcdef} (9) & Chấp nhận & Đăng ký OK & \textcolor{green!60!black}{Pass} \\
    TC-A-BV-04 & Ngoài biên dưới: rỗng & Form: để trống ô pw & Từ chối & Bị chặn (\texttt{required}+regex) & \textcolor{green!60!black}{Pass} \\
    TC-A-BV-05 & Biên trên khuyết: tên/MK rất dài & Form: dán $\approx$5.000 ký tự (MK vẫn hợp regex) & Nên giới hạn độ dài & Chấp nhận, không giới hạn & \textcolor{orange}{Note} \\
\end{longtable}
\noindentTC-A-BV-05 quan sát được trên form (dán chuỗi dài rồi Đăng ký) nhưng hệ quả chỉ
là ``vùng xám'' robustness (không có ràng buộc độ dài trên) --- ghi nhận là điểm cần cải thiện,
không phải pass/fail dứt khoát theo đặc tả.

\subsection{Bước 3: Ma Trận Phủ Biên (off / on / on$+1$)}
Bảng dưới chứng minh mỗi biên ở Bước~1 đều được kiểm đủ ba giá trị (kèm điểm ngoài miền), lấy ca
từ cả bảng BVA lẫn Domain Testing.

\begin{table}[H]
    \centering
    \caption{Phủ biên — Feature A (FR-01)}
    \begin{tabularx}{\linewidth}{L{2.8cm} C{1.7cm} C{1.7cm} C{1.7cm} X}
        \toprule
        \textbf{Biên} & \textbf{Off} & \textbf{On} & \textbf{On$+1$} & \textbf{Ngoài / Ghi chú} \\
        \midrule
        password.length $\ge 8$ & BV-01 (7) & BV-02 (8) & BV-03 (9) & Ngoài miền: BV-04 (0). Không có max $\Rightarrow$ robustness BV-05. \\
        name.length $\ge 1$ & UI \texttt{required} (0) & ca hợp lệ ($\ge 1$) & --- & Biên đơn điểm; form chặn rỗng; không có max $\Rightarrow$ BV-05. \\
        email.length $\ge 1$ & UI \texttt{required} (0) & ca hợp lệ ($\ge 1$) & --- & Tương tự; không có biên trên. \\
        \bottomrule
    \end{tabularx}
\end{table}
\noindent\textbf{Kết luận:} \texttt{password.length} phủ đủ off/on/on$+1$ + ngoài miền (form chặn
rỗng ở BV-04); \texttt{name}/\texttt{email} là biên đơn điểm (form \texttt{required} chặn
rỗng) --- phát hiện chính là thiếu biên trên, ghi nhận ở BV-05.

\section{AI Gap Analysis}
Sau khi AI sinh bộ ca kiểm thử theo phân hoạch/biên, ta rà soát lại để tìm những ca dễ bị bỏ
sót. Lý do bỏ sót được phân loại: chất lượng prompt, giới hạn công cụ AI, hoặc
độ phức tạp nội tại của tính năng.

\begin{table}[H]
    \centering
    \caption{AI Gap Analysis — Feature A (FR-01)}
    \begin{tabularx}{\linewidth}{C{0.5cm} X L{4cm} C{1.8cm}}
        \toprule
        \textbf{\#} & \textbf{Ca/bug bị bỏ sót} & \textbf{Lý do bỏ sót} & \textbf{Bổ sung ở TC} \\
        \midrule
        1 & Đảo ngược điều kiện regex: bắt buộc khoảng trắng và cấm ký tự đặc biệt. Đọc chữ trợ giúp dễ giả định ngược lại. & Giới hạn công cụ + chất lượng prompt (không đọc kỹ regex thực thi). & TC-A-DT-01/02/05 \\
        2 & Email trùng phá vỡ định danh khi đăng nhập (tương tác đa tính năng/trạng thái); đối chiếu qua tab Người dùng của admin. & Độ phức tạp nội tại (cross-feature register$\to$login). & TC-A-DT-07 \\
        3 & Email có khoảng trắng đầu/cuối tạo bản ghi ``trùng nhưng khác'' (không chuẩn hóa). & Độ phức tạp (chuẩn hóa/locale). & EC-E5 (bổ sung) \\
        4 & Không có giới hạn độ dài trên cho các trường (biên trên không xác định). & Độ phức tạp/robustness; BVA thường chỉ nhìn biên đã đặc tả. & TC-A-BV-05 \\
        \bottomrule
    \end{tabularx}
\end{table}

% ============================================================
%  CHƯƠNG 3: FEATURE B — POOL B
% ============================================================
\chapter{Feature B — Mã giảm giá (Pool B)}

\section{Mô Tả Tính Năng}
Tính năng FR-09 — Mã giảm giá cho phép người dùng nhập một mã (\texttt{code})
để giảm giá cho đơn hàng có tổng tiền \texttt{total\_amount}. Logic cốt lõi nằm ở API
\texttt{POST /api/apply-coupon} (\texttt{backend/server.js}); dữ liệu mã lưu ở bảng
\texttt{coupons} (\texttt{backend/database.js}).

\begin{itemize}
    \item \textbf{Bảng \texttt{coupons}:} \texttt{code} (UNIQUE), \texttt{type}
    (\texttt{'percent'}/\texttt{'fixed'}), \texttt{discount\_value} (INTEGER),
    \texttt{min\_order\_amount} (INTEGER), \texttt{expired\_at} (DATETIME),
    \texttt{is\_active} (0/1), \texttt{max\_uses\_per\_user} (INTEGER).
    \item \textbf{API \texttt{apply-coupon}:} nhận \verb|{code, total_amount, user_id}|.
    Tìm mã đang \texttt{is\_active=1}; nếu \texttt{total\_amount > min\_order\_amount} thì
    kiểm tra hạn dùng, kiểm tra số lần dùng (chỉ khi có \texttt{user\_id}), rồi tính
    \texttt{discount\_amount} và \texttt{final\_amount}.
    \item \textbf{Các mã mẫu (seed):} \texttt{SAVE10} (percent 10, min 300k),
    \texttt{BIGBUY} (fixed 50k, min 500k), \texttt{VIP100} (fixed 100k, min 300k, tối đa 2 lần),
    \texttt{EXPIRED} (percent 20, hết hạn 2020).
\end{itemize}

\noindent\textbf{Đầu vào (trên trang Checkout):} ô Tổng tiền thanh toán
(\texttt{total\_amount}, chỉnh được), ô Mã Giảm Giá (\texttt{code}); \texttt{user\_id}
suy từ trạng thái đăng nhập (guest $=$ vắng).\\
\textbf{Đầu ra quan sát trên UI:} dòng ``Tiết kiệm'' (\texttt{discount\_amount}), ``Thành tiền''
(\texttt{final\_amount}) hiện màu xanh, hoặc thông báo lỗi màu đỏ dưới ô mã.

\vspace{4pt}
\noindent\textbf{Công thức tính giảm giá theo mã nguồn:}
\begin{lstlisting}[caption={Tính discount trong apply-coupon (server.js)},numbers=none]
if (total_amount > coupon.min_order_amount) { ... }        // (A) dùng '>'
if (coupon.type === "percent")
    discount_amount = Math.floor(total_amount * (1 - coupon.discount_value)); // (B)
else
    discount_amount = coupon.discount_value;
final_amount = total_amount - discount_amount;
\end{lstlisting}

\begin{notebox}
Hai khiếm khuyết trọng tâm của chương này: (A) dùng toán tử \texttt{>} nên đơn hàng
bằng đúng \texttt{min\_order\_amount} bị từ chối (lỗi off-by-one biên). (B)
\texttt{discount\_value} của mã \texttt{percent} lưu dạng số nguyên (\texttt{10} nghĩa là 10\%),
nhưng công thức lại là \texttt{total\_amount $\times$ (1 $-$ discount\_value)}. Với
\texttt{SAVE10} (10\%) trên đơn 400.000\,₫: đáng lẽ giảm 40.000 (còn 360.000), nhưng mã tính
\texttt{floor(400000 $\times$ (1$-$10)) $= -3.600.000$}, khiến \texttt{final\_amount $=
400000 - (-3.600.000) = 4.000.000$} — đắt gấp 10 lần. Công thức chỉ đúng cho mã
\texttt{fixed}.
\end{notebox}

\begin{notebox}
\textbf{Phạm vi \& quy ước cột kết quả.} Theo giảng viên (forum Q\&A HW02): functional testing
từ UI Frontend. Mọi ca thao tác trên \textbf{trang Checkout của web} (thêm sản phẩm vào
giỏ, chỉnh ô Tổng tiền, nhập mã, bấm ``Áp dụng''); ô Tổng tiền cho phép nhập trực tiếp nên kiểm
được biên \texttt{min\_order\_amount} ngay trên giao diện. Ca cần một mã đặc biệt (vd mã fixed
lớn hơn đơn) được tạo trước ở tab ``Mã Giảm Giá'' của app admin rồi áp trên Checkout.
Cột Kết quả/Trạng thái là con số/thông báo hiển thị trên màn hình; chế độ
``sinh viên tự chạy, AI dự đoán''.
\end{notebox}

\section{Domain Testing}

Áp dụng Domain Testing theo bốn bước (xác định biến $\to$ phân hoạch miền
$\to$ chọn điểm ON/OFF/IN/OUT $\to$ thiết kế ca). Với FR-09, biến có thứ tự trọng tâm là
\texttt{total\_amount} (biên tại \texttt{min\_order\_amount}) và \texttt{max\_uses\_per\_user};
các biến rời rạc (\texttt{code}, \texttt{type}, \texttt{is\_active}, \texttt{expiry}) được phủ
bằng phân hoạch tương đương --- mỗi trạng thái một lớp. Biến đầu ra (\texttt{discount\_amount},
\texttt{final\_amount}) cũng được đưa vào để bắt lỗi công thức. Bước~5 tổng hợp ma trận truy vết
EC $\to$ TC.

\subsection{Bước 1: Xác Định Các Biến Đầu Vào}
\begin{table}[H]
    \centering
    \caption{Danh sách biến — Feature B (FR-09)}
    \begin{tabularx}{\linewidth}{C{0.5cm} L{3cm} L{2.6cm} C{1.8cm} X}
        \toprule
        \textbf{\#} & \textbf{Biến} & \textbf{Kiểu} & \textbf{Vào/Ra} & \textbf{Ghi chú} \\
        \midrule
        1 & code & string & Input & Bắt buộc; so khớp phân biệt hoa/thường (SQLite BINARY). \\
        2 & total\_amount & integer (₫) & Input & Biến có thứ tự $\Rightarrow$ biên tại \texttt{min\_order\_amount}. \\
        3 & user\_id & integer / vắng & Input & Quyết định có kiểm tra số lần dùng hay không. \\
        4 & coupon.type & enum & Dữ liệu mã & \texttt{percent} (công thức sai) vs \texttt{fixed} (đúng). \\
        5 & coupon.discount\_value & integer & Dữ liệu mã & percent: 10 = 10\%; fixed: số tiền ₫. \\
        6 & coupon.min\_order\_amount & integer & Dữ liệu mã & Biên (so bằng \texttt{>}). \\
        7 & coupon.expired\_at & date & Dữ liệu mã & Biên so với now. \\
        8 & coupon.max\_uses\_per\_user & integer & Dữ liệu mã & Biên số lần dùng (\texttt{>=}). \\
        9 & coupon.is\_active & boolean & Dữ liệu mã & \texttt{0} $\Rightarrow$ coi như không tồn tại. \\
        10 & discount\_amount / final\_amount & integer & \textbf{Output} & Biên đầu ra: không âm, final $\ge 0$. \\
        \bottomrule
    \end{tabularx}
\end{table}

\subsection{Bước 2: Phân Hoạch Miền}
Dấu * đánh dấu lớp mà một hệ thống đúng phải xử lý theo đặc tả nhưng mã hiện tại
xử lý sai.

\subsubsection{Biến code}
\begin{table}[H]
    \centering
    \caption{Phân hoạch miền — code}
    \begin{tabularx}{\linewidth}{C{1.4cm} L{6cm} C{1.6cm} X}
        \toprule
        \textbf{EC} & \textbf{Điều kiện} & \textbf{Loại} & \textbf{Ví dụ} \\
        \midrule
        EC-C1 & Mã tồn tại, đang \texttt{is\_active=1} & Valid & \texttt{SAVE10} \\
        EC-C2 & Mã không tồn tại & Invalid & \texttt{NOPE99} \\
        EC-C3 & Rỗng / thiếu trường & Invalid & \verb|""| \\
        EC-C4 & Mã tồn tại nhưng \texttt{is\_active=0} & Invalid & (mã bị vô hiệu) \\
        EC-C5 & Đúng mã nhưng khác hoa/thường & Invalid* & \texttt{save10} \\
        \bottomrule
    \end{tabularx}
\end{table}

\subsubsection{Biến total\_amount (so với \texttt{min\_order\_amount})}
\begin{table}[H]
    \centering
    \caption{Phân hoạch miền — total\_amount (mã VIP100, min = 300.000)}
    \begin{tabularx}{\linewidth}{C{1.4cm} L{6cm} C{1.6cm} X}
        \toprule
        \textbf{EC} & \textbf{Điều kiện} & \textbf{Loại} & \textbf{Ví dụ} \\
        \midrule
        EC-T1 & \texttt{total > min} & Valid & 300.001 \\
        EC-T2 & \texttt{total = min} (đúng biên) & Valid* & 300.000 \\
        EC-T3 & \texttt{total < min} & Invalid & 299.999 \\
        EC-T4 & Âm hoặc 0 & Invalid & \texttt{-100} \\
        EC-T5 & Thiếu / không phải số & Invalid & \texttt{undefined} \\
        \bottomrule
    \end{tabularx}
\end{table}

\subsubsection{Biến user\_id \& loại mã}
\begin{table}[H]
    \centering
    \caption{Phân hoạch miền — user\_id, type, expiry}
    \begin{tabularx}{\linewidth}{C{1.4cm} L{6.5cm} C{1.6cm} X}
        \toprule
        \textbf{EC} & \textbf{Điều kiện} & \textbf{Loại} & \textbf{Ví dụ} \\
        \midrule
        EC-U1 & Có \texttt{user\_id}, chưa đạt \texttt{max\_uses} & Valid & dùng lần 1/2 \\
        EC-U2 & Có \texttt{user\_id}, đã đạt \texttt{max\_uses} & Invalid & dùng lần 3 (VIP100) \\
        EC-U3 & Vắng \texttt{user\_id} $\Rightarrow$ bỏ qua kiểm tra & Invalid* & không gửi \texttt{user\_id} \\
        EC-Y1 & Loại \texttt{fixed} & Valid & VIP100/BIGBUY \\
        EC-Y2 & Loại \texttt{percent} (công thức sai) & Valid* & SAVE10 \\
        EC-X1 & Còn hạn & Valid & \texttt{2099-12-31} \\
        EC-X2 & Hết hạn & Invalid & EXPIRED (2020) \\
        \bottomrule
    \end{tabularx}
\end{table}

\subsection{Bước 3: Chọn Điểm Kiểm Thử (ON/OFF/IN/OUT)}
Biên số học rõ ràng nhất là \texttt{total\_amount} tại \texttt{min\_order\_amount}. Dùng mã
VIP100 (\texttt{fixed} 100.000, min 300.000) để cô lập biên chấp nhận/từ chối, tránh
để lỗi công thức percent làm nhiễu.

\begin{table}[H]
    \centering
    \caption{Điểm ON/OFF/IN/OUT — total\_amount (min = 300.000)}
    \begin{tabularx}{\linewidth}{L{2.4cm} C{2.4cm} X}
        \toprule
        \textbf{Loại điểm} & \textbf{Giá trị} & \textbf{Kỳ vọng (đặc tả)} \\
        \midrule
        ON-point  & 300.000 & Chấp nhận (đơn đủ mức tối thiểu) \\
        OFF-point & 299.999 & Từ chối (dưới mức tối thiểu) \\
        IN-point  & 400.000 & Chấp nhận \\
        OUT-point & 0       & Từ chối \\
        \bottomrule
    \end{tabularx}
\end{table}

\noindent\textbf{Nhận xét:} do mã dùng \texttt{total\_amount > min} (toán tử \texttt{>}),
ON-point 300.000 bị từ chối trong khi đặc tả yêu cầu chấp nhận — đây là
lỗi off-by-one biên (BUG-B-02). Biên \texttt{max\_uses\_per\_user} lại dùng \texttt{>=} nên
được cài đặt đúng (sẽ Pass), ngoại trừ lỗ hổng bỏ qua khi vắng \texttt{user\_id}.

\subsection{Bước 4: Thiết Kế Test Cases — Domain Testing}
\begin{longtable}{C{1.1cm} L{3.1cm} L{3.2cm} L{3cm} C{2cm} C{1.4cm}}
    \caption{Test cases — Domain Testing, Feature B} \\
    \toprule
    \textbf{TC\#} & \textbf{Mô tả} & \textbf{Input} & \textbf{Expected Output} & \textbf{Kết quả} & \textbf{Trạng thái} \\
    \midrule
    \endfirsthead
    \multicolumn{6}{c}{\small\itshape (tiếp theo)} \\
    \toprule
    \textbf{TC\#} & \textbf{Mô tả} & \textbf{Input} & \textbf{Expected Output} & \textbf{Kết quả} & \textbf{Trạng thái} \\
    \midrule
    \endhead
    \midrule
    \multicolumn{6}{r}{\small\itshape (còn tiếp)} \\
    \endfoot
    \bottomrule
    \endlastfoot
    TC-B-DT-01 & Mã percent hợp lệ \newline (BUG-B-01) & Checkout: tổng 400.000, nhập SAVE10, Áp dụng & Tiết kiệm 40.000, Thành tiền 360.000 & Tiết kiệm $-3.600.000$, Thành tiền 4.000.000 & \textcolor{red}{Fail} \\
    TC-B-DT-02 & Mã fixed hợp lệ & Checkout: tổng 600.000, mã BIGBUY & Tiết kiệm 50.000, Thành tiền 550.000 & Như expected & \textcolor{green!60!black}{Pass} \\
    TC-B-DT-03 & Mã fixed hợp lệ & Checkout: tổng 400.000, mã VIP100 & Tiết kiệm 100.000, Thành tiền 300.000 & Như expected & \textcolor{green!60!black}{Pass} \\
    TC-B-DT-04 & Mã không tồn tại & Checkout: nhập NOPE99 & Báo lỗi ``không tồn tại'' & Báo lỗi ``không tồn tại'' & \textcolor{green!60!black}{Pass} \\
    TC-B-DT-05 & Mã hết hạn & Checkout: tổng 200.000, mã EXPIRED & Báo lỗi ``đã hết hạn'' & Báo lỗi ``đã hết hạn'' & \textcolor{green!60!black}{Pass} \\
    TC-B-DT-06 & total = min (đúng biên) \newline (BUG-B-02) & Checkout: tổng 300.000, mã VIP100 & Chấp nhận (Thành tiền 200.000) & Báo ``chưa đủ giá trị'' & \textcolor{red}{Fail} \\
    TC-B-DT-07 & Đạt giới hạn số lần dùng & Đăng nhập, checkout VIP100 đủ 2 lần, áp lần 3 & Báo ``đã đạt giới hạn'' & Báo ``đã đạt giới hạn'' & \textcolor{green!60!black}{Pass} \\
    TC-B-DT-08 & Khách (guest) không bị giới hạn số lần \newline (BUG-B-03) & Không đăng nhập, áp VIP100 nhiều lần ở Checkout & Vẫn phải chặn theo giới hạn & Luôn áp dụng được (guest bỏ kiểm tra) & \textcolor{orange}{Note} \\
    TC-B-DT-09 & fixed vượt tổng đơn $\Rightarrow$ âm \newline (BUG-B-04) & Tạo mã fixed (discount 100.000, min 50.000) ở admin; Checkout tổng 60.000, áp mã & Thành tiền $\ge 0$, có chặn & Thành tiền $= -40.000$ (không chặn) & \textcolor{red}{Fail} \\
\end{longtable}

\subsection{Bước 5: Ma Trận Truy Vết Độ Phủ (EC $\to$ Test Case)}
Bảng dưới đối chiếu mọi lớp tương đương (Bước~2) với ca kiểm thử phủ nó (kể cả ca BVA). Các lớp
cùng nhánh xử lý dùng chung ca đại diện; một lớp chưa có ca riêng được ghi nhận minh bạch để bổ
sung.

\begin{table}[H]
    \centering
    \caption{Truy vết EC $\to$ TC — Feature B (FR-09)}
    \begin{tabularx}{\linewidth}{C{1.8cm} C{3.4cm} X}
        \toprule
        \textbf{EC} & \textbf{Ca phủ (DT/BVA)} & \textbf{Ghi chú} \\
        \midrule
        EC-C1 & TC-B-DT-01/02/03 & Mã tồn tại, đang \texttt{is\_active=1}. \\
        EC-C2 & TC-B-DT-04 & Mã không tồn tại. \\
        EC-C3 & (UI: nút Áp dụng bị khóa) & Ô mã rỗng thì nút ``Áp dụng'' bị vô hiệu $\Rightarrow$ không gửi được --- hành vi đúng. \\
        EC-C4 & TC-B-DT-04 (đại diện) & Truy vấn lọc \texttt{is\_active=1} nên mã inactive đi cùng nhánh ``không tồn tại''; không có mã seed inactive. \\
        EC-C5 & (UI ép chữ hoa) & Ô mã tự chuyển hoa (\texttt{toUpperCase}) $\Rightarrow$ không gửi được chữ thường; lỗi phân biệt hoa/thường của backend bị UI che. \\
        EC-T1 & TC-B-DT-02/03, BV-03 & \texttt{total > min}. \\
        EC-T2 & TC-B-DT-06, BV-02 & \texttt{total = min} (biên off-by-one). \\
        EC-T3 & TC-B-BV-01 & \texttt{total < min}. \\
        EC-T4 & TC-B-BV-01 (đại diện) & 0/âm $<$ min $\Rightarrow$ cùng nhánh từ chối như EC-T3 (OUT-point ở Bước~3). \\
        EC-T5 & --- & Ô Tổng tiền là \texttt{type="number"} nên UI chỉ nhập được số; lớp ``không phải số'' không đạt được từ giao diện. \\
        EC-U1 & TC-B-DT-03, BV-04 & Có \texttt{user\_id} (đăng nhập), chưa đạt max. \\
        EC-U2 & TC-B-DT-07, BV-05 & Có \texttt{user\_id}, đã đạt max. \\
        EC-U3 & TC-B-DT-08 & Guest (vắng \texttt{user\_id}) $\Rightarrow$ bỏ kiểm tra. \\
        EC-Y1 & TC-B-DT-02/03 & Loại \texttt{fixed}. \\
        EC-Y2 & TC-B-DT-01 & Loại \texttt{percent} (công thức sai). \\
        EC-X1 & TC-B-BV-06 & Còn hạn. \\
        EC-X2 & TC-B-DT-05, BV-06 & Hết hạn. \\
        \bottomrule
    \end{tabularx}
\end{table}
\noindent\textbf{Kết luận độ phủ:} mọi lớp đều có ca phủ hoặc được UI chặn ngay (mã rỗng, chữ
thường); riêng EC-T5 (Tổng tiền không phải số) không đạt được từ giao diện vì ô là kiểu số.

\section{Boundary Value Analysis (BVA)}

BVA cho FR-09 nhắm vào ba biên: \texttt{total\_amount} so với \texttt{min\_order\_amount},
\texttt{max\_uses\_per\_user}, và \texttt{expired\_at} so với now. Với mỗi biên số, ta sinh
bộ ba \texttt{off/on/on$+1$} để lộ off-by-one --- điển hình là toán tử \texttt{>} tại biên
\texttt{min} khiến đơn bằng đúng mức tối thiểu bị từ chối. Bước~3 tổng hợp ma trận phủ biên, cho
thấy mỗi biên đều có đủ điểm dưới/trên.

\subsection{Bước 1: Xác Định Ranh Giới}
\begin{table}[H]
    \centering
    \caption{Bảng giá trị biên — Feature B (FR-09)}
    \begin{tabularx}{\linewidth}{L{3.4cm} C{1.8cm} C{1.8cm} C{1.8cm} X}
        \toprule
        \textbf{Biên} & \textbf{Off} & \textbf{On} & \textbf{On+1} & \textbf{Ghi chú} \\
        \midrule
        total\_amount vs min (VIP100=300k) & 299.999 & 300.000 & 300.001 & On bị từ chối (dùng \texttt{>}) $\Rightarrow$ lỗi. \\
        max\_uses\_per\_user (VIP100=2) & 1 lần & 2 lần & 3 lần & Dùng \texttt{>=}, cài đặt đúng. \\
        expired\_at vs now & quá khứ & --- & tương lai & So sánh \texttt{expiry < now} đúng. \\
        \bottomrule
    \end{tabularx}
\end{table}

\subsection{Bước 2: Thiết Kế Test Cases — BVA}
\begin{longtable}{C{1.1cm} L{3.2cm} L{3.2cm} L{3cm} C{2cm} C{1.4cm}}
    \caption{Test cases — BVA, Feature B} \\
    \toprule
    \textbf{TC\#} & \textbf{Mô tả} & \textbf{Input} & \textbf{Expected Output} & \textbf{Kết quả} & \textbf{Trạng thái} \\
    \midrule
    \endfirsthead
    \toprule
    \textbf{TC\#} & \textbf{Mô tả} & \textbf{Input} & \textbf{Expected Output} & \textbf{Kết quả} & \textbf{Trạng thái} \\
    \midrule
    \endhead
    \bottomrule
    \endlastfoot
    TC-B-BV-01 & min$-1$: dưới biên & Checkout: tổng 299.999, VIP100 & Từ chối (chưa đủ) & Từ chối & \textcolor{green!60!black}{Pass} \\
    TC-B-BV-02 & min: đúng biên \newline (BUG-B-02) & Checkout: tổng 300.000, VIP100 & Chấp nhận, Thành tiền 200.000 & Báo ``chưa đủ'' & \textcolor{red}{Fail} \\
    TC-B-BV-03 & min$+1$: trên biên & Checkout: tổng 300.001, VIP100 & Chấp nhận, Thành tiền 200.001 & Chấp nhận & \textcolor{green!60!black}{Pass} \\
    TC-B-BV-04 & max$-1$ lần dùng & Đăng nhập, đã checkout VIP100 1 lần, áp lần 2 & Chấp nhận & Chấp nhận & \textcolor{green!60!black}{Pass} \\
    TC-B-BV-05 & max lần dùng (đúng biên) & Đăng nhập, đã checkout VIP100 2 lần, áp lần 3 & Báo ``đạt giới hạn'' & Từ chối & \textcolor{green!60!black}{Pass} \\
    TC-B-BV-06 & Biên hạn dùng & Checkout: EXPIRED (quá khứ) vs SAVE10 (tương lai) & Quá khứ từ chối, tương lai nhận & Đúng như vậy & \textcolor{green!60!black}{Pass} \\
\end{longtable}

\subsection{Bước 3: Ma Trận Phủ Biên (off / on / on$+1$)}
Bảng dưới chứng minh mỗi biên ở Bước~1 đều được kiểm đủ ba giá trị, lấy ca từ cả bảng BVA lẫn
Domain Testing.

\begin{table}[H]
    \centering
    \caption{Phủ biên — Feature B (FR-09)}
    \begin{tabularx}{\linewidth}{L{2.8cm} C{1.9cm} C{1.7cm} C{1.6cm} X}
        \toprule
        \textbf{Biên} & \textbf{Off} & \textbf{On} & \textbf{On$+1$} & \textbf{Ngoài / Ghi chú} \\
        \midrule
        total\_amount vs min (300k) & BV-01 (299.999) & BV-02 (300.000) & BV-03 (300.001) & Ngoài miền: 0/âm (OUT-point Bước~3). Off-by-one tại On (dùng \texttt{>}). \\
        max\_uses (2) & BV-04 (1) & BV-05 (2) & DT-07 (3) & On$+1$ = lần dùng thứ 3 bị từ chối (đúng, dùng \texttt{>=}). \\
        expired\_at vs now & DT-05/BV-06 (quá khứ) & now & BV-06 (tương lai) & Biên thời gian: quá khứ từ chối, tương lai nhận. \\
        \bottomrule
    \end{tabularx}
\end{table}
\noindent\textbf{Kết luận:} cả ba biên đều có điểm dưới/trên; biên \texttt{min\_order\_amount}
lộ off-by-one (BUG-B-02), còn \texttt{max\_uses} và \texttt{expiry} được cài đặt đúng.

\section{AI Gap Analysis}
\begin{table}[H]
    \centering
    \caption{AI Gap Analysis — Feature B (FR-09)}
    \begin{tabularx}{\linewidth}{C{0.5cm} X L{4cm} C{1.8cm}}
        \toprule
        \textbf{\#} & \textbf{Ca/bug bị bỏ sót} & \textbf{Lý do bỏ sót} & \textbf{Bổ sung ở TC} \\
        \midrule
        1 & Công thức percent sai (\texttt{1$-$discount\_value}) làm giá tăng gấp $\sim$10 lần — cần đọc phép tính thực tế, không chỉ tên biến. & Giới hạn công cụ + độ phức tạp (biến đầu ra, không phải input). & TC-B-DT-01 \\
        2 & Off-by-one biên \texttt{min\_order\_amount} (\texttt{>} thay vì \texttt{>=}) — chỉ lộ khi thử đúng giá trị biên. & Chất lượng prompt (không yêu cầu thử ca ``bằng đúng biên''). & TC-B-DT-06, BV-02 \\
        3 & Khách (guest, vắng \texttt{user\_id}) không bị giới hạn số lần dùng (nhánh else). & Độ phức tạp nội tại (phụ thuộc luồng điều khiển/nhánh). & TC-B-DT-08 \\
        4 & Không chặn \texttt{final\_amount} âm với mã fixed (tạo mã ở admin, áp trên Checkout). & Độ phức tạp (biên đầu ra + cần tự tạo mã qua UI admin). & TC-B-DT-09 \\
        \bottomrule
    \end{tabularx}
\end{table}

% ============================================================
%  CHƯƠNG 4: FEATURE C — POOL C
% ============================================================
\chapter{Feature C — Quản lý danh mục (Pool C — Web Admin)}

\section{Mô Tả Tính Năng}
Tính năng FR-14 — Quản lý danh mục (CRUD) thuộc trang Web Admin, cho phép
thêm/sửa/xóa danh mục sản phẩm. Toàn bộ logic nằm ở 4 endpoint trong
\texttt{backend/server.js}; bảng \texttt{categories} khai báo ở \texttt{backend/database.js}.

\begin{notebox}
\textbf{Phạm vi kiểm thử — kiểm từ UI trang Admin.} Theo giảng viên Hồ Tuấn Thanh (diễn đàn
Q\&A HW02, Moodle 01/07/2026): ``Mình thực hiện functional testing thì test từ UI
Frontend nhé''. FR-14 được thao tác qua ứng dụng admin riêng \texttt{frontend-admin}
(React/Vite), tab ``Danh mục'' (\texttt{frontend-admin/src/App.jsx}). Đăng nhập bằng
tài khoản admin có sẵn (\texttt{admin@eshop.com} / \texttt{Admin123!}); form login chặn tài
khoản không phải admin (``Bạn không phải là admin!''). Giao diện danh mục chỉ
cung cấp hai thao tác: Thêm mới (\texttt{Thêm mới}) và Xóa (\texttt{Xóa}) --- không có nút Sửa.
Vì vậy mọi kỳ vọng và cột Kết quả dưới đây là những gì quan sát được trên UI (bảng danh
mục, thông báo, hệ quả ở tab Sản phẩm), không phải gọi API trực tiếp.
\end{notebox}

\noindent Về mặt cài đặt, tab Danh mục gọi tới 4 endpoint sau (UI chỉ dùng POST tạo và DELETE
xóa; GET để nạp bảng); liệt kê để đối chiếu hành vi quan sát được:
\begin{itemize}
    \item \texttt{GET /api/categories} — công khai (không xác thực), trả về mọi danh mục.
    \item \texttt{POST /api/categories} — \texttt{authenticateToken}, body \verb|{name}| $\to$
    \texttt{INSERT INTO categories (name)}.
    \item \texttt{PUT /api/categories/:id} — \texttt{authenticateToken}, \verb|{name}| $\to$
    \texttt{UPDATE categories SET name=? WHERE id=?}.
    \item \texttt{DELETE /api/categories/:id} — \texttt{authenticateToken} $\to$
    \texttt{DELETE FROM categories WHERE id=?}.
\end{itemize}

\noindent\textbf{Lược đồ:} \texttt{categories(id INTEGER PK AUTOINCREMENT, name TEXT)} — cột
\texttt{name} không \texttt{NOT NULL}, không \texttt{UNIQUE}, không giới hạn độ
dài. Bảng \texttt{products} có \texttt{category\_id INTEGER} nhưng không có khóa ngoại
(foreign key) trỏ tới \texttt{categories}. \textbf{Seed:} 3 danh mục — \texttt{1=Điện thoại}
(có 2 sản phẩm), \texttt{2=Laptop} (1 sản phẩm), \texttt{3=Phụ kiện} (2 sản phẩm).

\vspace{4pt}
\noindent\textbf{Middleware \texttt{authenticateToken}} chỉ kiểm tra JWT hợp lệ
(\texttt{jwt.verify}) rồi gán \texttt{req.user}; không kiểm tra \texttt{role ===
"admin"}. Do đó bất kỳ người dùng đã đăng nhập (kể cả \texttt{role='user'}) đều gọi được
các endpoint quản trị này.

\begin{notebox}
Khiếm khuyết trọng tâm của chương: (1) thiếu kiểm soát truy cập theo vai trò — CRUD
danh mục (chức năng admin) mở cho mọi tài khoản đăng nhập; (2) không kiểm tra dữ liệu
\texttt{name} (rỗng/trùng/quá dài); (3) PUT/DELETE trên \texttt{id} không tồn tại vẫn
trả \texttt{200} thành công (không kiểm \texttt{this.changes}); (4) xóa danh mục làm
mồ côi sản phẩm (không khóa ngoại/cascade).
\end{notebox}

\begin{notebox}
\textbf{Quy ước cột kết quả} (như Feature A/B): mọi thao tác thực hiện trên giao diện app
admin (\texttt{frontend-admin}, tab ``Danh mục''); cột Kết quả/Trạng thái là hành
vi \textbf{quan sát được trên UI} --- bảng danh mục, thông báo (\texttt{alert}), và hệ quả nhìn
thấy ở tab ``Sản phẩm''. Chế độ ``sinh viên tự chạy, AI dự đoán từ mã nguồn''; sinh viên đăng
nhập \texttt{admin@eshop.com}/\texttt{Admin123!}, mở tab Danh mục và đối chiếu.
\end{notebox}

\section{Domain Testing}

Áp dụng Domain Testing theo bốn bước (xác định biến $\to$ phân hoạch miền
$\to$ chọn điểm ON/OFF/IN/OUT $\to$ thiết kế ca). Trên UI admin, tab Danh mục chỉ có hai thao
tác: Thêm mới (ô nhập ``Tên danh mục mới'') và Xóa (nút trên mỗi hàng) ---
không có chức năng Sửa. Do đó biến có thứ tự để áp dụng biên là \texttt{name.length}
(ô nhập không ràng buộc độ dài); các biến phân loại là tính hợp lệ của tên (rỗng / trùng /
ký tự đặc biệt) và trạng thái danh mục khi xóa (có / không có sản phẩm). Việc đăng nhập
admin là tiền điều kiện (form login chặn tài khoản không phải admin). Bước~5 tổng hợp ma
trận truy vết EC $\to$ TC.

\subsection{Bước 1: Xác Định Các Biến Đầu Vào}
\begin{table}[H]
    \centering
    \caption{Danh sách biến — Feature C (FR-14), nhìn từ UI admin}
    \begin{tabularx}{\linewidth}{C{0.5cm} L{3.2cm} L{2.6cm} C{1.8cm} X}
        \toprule
        \textbf{\#} & \textbf{Biến} & \textbf{Kiểu} & \textbf{Vào/Ra} & \textbf{Ghi chú} \\
        \midrule
        1 & name (ô ``Tên danh mục mới'') & string & Input & Form Thêm mới; không kiểm rỗng/trùng/độ dài. \\
        2 & name.length & integer (ẩn) & Input & Ô nhập không có min ($>0$)/max $\Rightarrow$ dùng BVA. \\
        3 & name.isUnique & boolean (ẩn) & Input & Không có \texttt{UNIQUE} $\Rightarrow$ cho trùng. \\
        4 & category.hasProducts & boolean (trạng thái) & Input & Ảnh hưởng toàn vẹn khi bấm Xóa. \\
        5 & đăng nhập admin & boolean (gate) & Tiền ĐK & Non-admin bị chặn (``Bạn không phải là admin!''). \\
        6 & phản hồi UI & bảng/thông báo & \textbf{Output} & Hàng trong bảng danh mục, \texttt{alert}, hệ quả tab Sản phẩm. \\
        \bottomrule
    \end{tabularx}
\end{table}

\subsection{Bước 2: Phân Hoạch Miền}
Dấu * = lớp mà hệ thống đúng phải xử lý theo đặc tả nhưng UI/mã hiện tại làm sai.

\subsubsection{Biến name (ô ``Tên danh mục mới'')}
\begin{table}[H]
    \centering
    \caption{Phân hoạch miền — name}
    \begin{tabularx}{\linewidth}{C{1.4cm} L{6cm} C{1.6cm} X}
        \toprule
        \textbf{EC} & \textbf{Điều kiện} & \textbf{Loại} & \textbf{Ví dụ (nhập trên UI)} \\
        \midrule
        EC-NM1 & Tên mới, không rỗng, chưa trùng & Valid & \texttt{Sách} \\
        EC-NM2 & Rỗng (để trống ô rồi bấm Thêm) & Invalid* & \verb|""| \\
        EC-NM3 & Chỉ khoảng trắng & Invalid* & \verb|"   "| \\
        EC-NM4 & Trùng tên đã tồn tại & Invalid* & \texttt{Laptop} \\
        EC-NM5 & Rất dài (ô không giới hạn) & Valid* & 5.000 ký tự \\
        EC-NM6 & Ký tự đặc biệt/injection & Valid (literal) & \texttt{x'; DROP TABLE categories;-{}-} \\
        \bottomrule
    \end{tabularx}
\end{table}

\subsubsection{Trạng thái danh mục khi Xóa}
\begin{table}[H]
    \centering
    \caption{Phân hoạch miền — trạng thái khi bấm Xóa}
    \begin{tabularx}{\linewidth}{C{1.4cm} L{6cm} C{1.6cm} X}
        \toprule
        \textbf{EC} & \textbf{Điều kiện} & \textbf{Loại} & \textbf{Kỳ vọng} \\
        \midrule
        EC-DL1 & Danh mục không có sản phẩm & Valid & Xóa được, biến mất khỏi bảng \\
        EC-DL2 & Danh mục đang có sản phẩm & Invalid* & Nên chặn/cảnh báo; tránh mồ côi \\
        \bottomrule
    \end{tabularx}
\end{table}

\subsubsection{Tiền điều kiện — quyền truy cập (form login admin)}
\begin{table}[H]
    \centering
    \caption{Phân hoạch miền — đăng nhập app admin}
    \begin{tabularx}{\linewidth}{C{1.4cm} L{6.5cm} C{1.6cm} X}
        \toprule
        \textbf{EC} & \textbf{Điều kiện} & \textbf{Loại} & \textbf{Kỳ vọng (quan sát UI)} \\
        \midrule
        EC-AU1 & Đăng nhập tài khoản admin & Valid & Vào được, thấy tab Danh mục \\
        EC-AU2 & Đăng nhập tài khoản thường & Invalid & Bị chặn: ``Bạn không phải là admin!'' \\
        \bottomrule
    \end{tabularx}
\end{table}
\noindentGhi chú: từ UI, tài khoản thường không vào được trang admin nên không thể
thao tác CRUD danh mục — kiểm soát truy cập ở tầng giao diện hoạt động đúng. (Việc backend không
kiểm \texttt{role} khi gọi API trực tiếp nằm ngoài phạm vi functional testing từ UI.)

\subsection{Bước 3: Chọn Điểm Kiểm Thử (ON/OFF/IN/OUT)}
Biến có thứ tự thao tác được từ UI là \texttt{name.length} của ô ``Tên danh mục mới''. Theo tinh
thần đặc tả, tên danh mục hợp lệ phải không rỗng (độ dài $\ge 1$); ô nhập lại không có
giới hạn trên. Áp dụng ON/OFF/IN trên biên rỗng.

\begin{table}[H]
    \centering
    \caption{Điểm ON/OFF/IN/OUT — name.length (biên hợp lệ $\ge 1$)}
    \begin{tabularx}{\linewidth}{L{2.4cm} C{2.6cm} X}
        \toprule
        \textbf{Loại điểm} & \textbf{name.length} & \textbf{Kỳ vọng (đặc tả)} \\
        \midrule
        OFF (ngoài) & 0 (rỗng) & Từ chối, không tạo danh mục \\
        ON (biên)   & 1 & Chấp nhận (\texttt{"A"}) \\
        IN (trong)  & 4 & Chấp nhận (\texttt{"Sách"}) \\
        (biên trên) & không có max & Nên có giới hạn hợp lý; UI để vô hạn \\
        \bottomrule
    \end{tabularx}
\end{table}

\noindent\textbf{Nhận xét:} ô nhập và backend đều không kiểm rỗng, nên điểm OFF (độ dài 0)
vẫn tạo được một danh mục tên trống hiển thị trong bảng --- lỗi (BUG-C-01); đồng thời không có
biên trên nên tên cực dài vẫn được nhận (BUG-C-04).

\subsection{Bước 4: Thiết Kế Test Cases — Domain Testing}
\begin{longtable}{C{1.2cm} L{3.1cm} L{3.1cm} L{3cm} C{1.9cm} C{1.4cm}}
    \caption{Test cases — Domain Testing, Feature C} \\
    \toprule
    \textbf{TC\#} & \textbf{Mô tả} & \textbf{Input} & \textbf{Expected Output} & \textbf{Kết quả} & \textbf{Trạng thái} \\
    \midrule
    \endfirsthead
    \multicolumn{6}{c}{\small\itshape (tiếp theo)} \\
    \toprule
    \textbf{TC\#} & \textbf{Mô tả} & \textbf{Input} & \textbf{Expected Output} & \textbf{Kết quả} & \textbf{Trạng thái} \\
    \midrule
    \endhead
    \midrule
    \multicolumn{6}{r}{\small\itshape (còn tiếp)} \\
    \endfoot
    \bottomrule
    \endlastfoot
    TC-C-DT-01 & Thêm danh mục hợp lệ & Ô tên = ``Sách'', bấm Thêm mới & Hàng ``Sách'' xuất hiện trong bảng & Hàng ``Sách'' hiện ra & \textcolor{green!60!black}{Pass} \\
    TC-C-DT-02 & Thêm tên rỗng \newline (BUG-C-01) & Để trống ô, bấm Thêm mới & Chặn / báo lỗi, không tạo & Tạo 1 hàng tên trống & \textcolor{red}{Fail} \\
    TC-C-DT-03 & Thêm tên chỉ khoảng trắng \newline (BUG-C-01) & Ô tên = \verb|"   "| & Chặn / báo lỗi & Tạo 1 hàng tên trống & \textcolor{red}{Fail} \\
    TC-C-DT-04 & Thêm tên trùng \newline (BUG-C-02) & Ô tên = ``Laptop'' (đã có) & Từ chối trùng & Tạo ``Laptop'' thứ 2 & \textcolor{red}{Fail} \\
    TC-C-DT-05 & Đăng nhập tài khoản thường vào app admin & Login \texttt{user@\ldots} & Bị chặn, không vào & ``Bạn không phải là admin!'' & \textcolor{green!60!black}{Pass} \\
    TC-C-DT-06 & Thêm tên có ký tự đặc biệt/injection & Ô tên = \texttt{x'; DROP TABLE\ldots} & Lưu literal an toàn, hiện đúng chuỗi & Hiện đúng literal, DB an toàn & \textcolor{green!60!black}{Pass} \\
    TC-C-DT-07 & Xóa danh mục không có sản phẩm & Bấm Xóa ở danh mục vừa thêm & Biến mất khỏi bảng & Biến mất khỏi bảng & \textcolor{green!60!black}{Pass} \\
    TC-C-DT-08 & Xóa danh mục đang có sản phẩm \newline (BUG-C-03) & Bấm Xóa ở ``Điện thoại'' (2 sp) & Chặn/cảnh báo, tránh mồ côi & Xóa ngay; tab Sản phẩm còn SP mồ côi & \textcolor{red}{Fail} \\
\end{longtable}

\subsection{Bước 5: Ma Trận Truy Vết Độ Phủ (EC $\to$ Test Case)}
Bảng dưới đối chiếu mọi lớp tương đương (Bước~2) với ca kiểm thử phủ nó (kể cả ca BVA). Lớp cùng
kỳ vọng hành vi dùng chung ca đại diện theo nguyên tắc phân hoạch tương đương.

\begin{table}[H]
    \centering
    \caption{Truy vết EC $\to$ TC — Feature C (FR-14)}
    \begin{tabularx}{\linewidth}{C{1.8cm} C{3.4cm} X}
        \toprule
        \textbf{EC} & \textbf{Ca phủ (DT/BVA)} & \textbf{Ghi chú} \\
        \midrule
        EC-NM1 & TC-C-DT-01, BV-02 & Tên mới hợp lệ. \\
        EC-NM2 & TC-C-DT-02, BV-01 & Tên rỗng (độ dài 0). \\
        EC-NM3 & TC-C-DT-03 & Chỉ khoảng trắng (cùng kỳ vọng từ chối). \\
        EC-NM4 & TC-C-DT-04 & Trùng tên đã tồn tại. \\
        EC-NM5 & TC-C-BV-03 & Tên cực dài (biên trên khuyết). \\
        EC-NM6 & TC-C-DT-06 & Ký tự đặc biệt/injection. \\
        EC-DL1 & TC-C-DT-07 & Xóa danh mục không có sản phẩm. \\
        EC-DL2 & TC-C-DT-08 & Xóa danh mục đang có sản phẩm. \\
        EC-AU1 & TC-C-DT-01 & Đăng nhập admin (vào được). \\
        EC-AU2 & TC-C-DT-05 & Đăng nhập tài khoản thường (bị chặn). \\
        \bottomrule
    \end{tabularx}
\end{table}
\noindent\textbf{Kết luận độ phủ:} 6/6 lớp \texttt{name}, 2/2 lớp trạng thái khi xóa, 2/2 lớp
quyền truy cập đều có ít nhất một ca kiểm thử.

\section{Boundary Value Analysis (BVA)}

BVA cho FR-14 (từ UI) nhắm vào biên độ dài tên nhập ở ô ``Tên danh mục mới'':
\texttt{name.length}. Ta kiểm biên rỗng (0/1) và biên trên (ô không có max). Đây là biên duy
nhất thao tác được từ giao diện (không có nút Sửa, không nhập được \texttt{id} tùy ý). Bước~3
tổng hợp ma trận phủ biên.

\subsection{Bước 1: Xác Định Ranh Giới}
\begin{table}[H]
    \centering
    \caption{Bảng giá trị biên — Feature C (FR-14)}
    \begin{tabularx}{\linewidth}{L{3.2cm} C{1.8cm} C{1.8cm} C{2cm} X}
        \toprule
        \textbf{Biên} & \textbf{Off/dưới} & \textbf{On} & \textbf{Off/trên} & \textbf{Ghi chú} \\
        \midrule
        name.length & 0 (rỗng) & 1 & --- (không max) & Độ dài 0 vẫn tạo được (BUG-C-01); không có biên trên (BUG-C-04). \\
        \bottomrule
    \end{tabularx}
\end{table}

\subsection{Bước 2: Thiết Kế Test Cases — BVA}
\begin{longtable}{C{1.2cm} L{3.2cm} L{3.1cm} L{3cm} C{1.9cm} C{1.4cm}}
    \caption{Test cases — BVA, Feature C} \\
    \toprule
    \textbf{TC\#} & \textbf{Mô tả} & \textbf{Input (UI)} & \textbf{Expected Output} & \textbf{Kết quả} & \textbf{Trạng thái} \\
    \midrule
    \endfirsthead
    \toprule
    \textbf{TC\#} & \textbf{Mô tả} & \textbf{Input (UI)} & \textbf{Expected Output} & \textbf{Kết quả} & \textbf{Trạng thái} \\
    \midrule
    \endhead
    \bottomrule
    \endlastfoot
    TC-C-BV-01 & name rỗng (length 0) \newline (BUG-C-01) & Để trống ô, bấm Thêm & Từ chối, không tạo & Tạo hàng tên trống & \textcolor{red}{Fail} \\
    TC-C-BV-02 & name length 1 & Ô tên = ``A'' & Thêm được & Hàng ``A'' hiện ra & \textcolor{green!60!black}{Pass} \\
    TC-C-BV-03 & name cực dài \newline (BUG-C-04) & Ô tên = 5.000 ký tự & Nên giới hạn/từ chối & Tạo được, lưu toàn bộ & \textcolor{red}{Fail} \\
\end{longtable}

\subsection{Bước 3: Ma Trận Phủ Biên (off / on / on$+1$)}
Bảng dưới chứng minh biên ở Bước~1 được kiểm đủ.

\begin{table}[H]
    \centering
    \caption{Phủ biên — Feature C (FR-14)}
    \begin{tabularx}{\linewidth}{L{2.8cm} C{1.7cm} C{2.4cm} C{2cm} X}
        \toprule
        \textbf{Biên} & \textbf{Dưới} & \textbf{On} & \textbf{Trên} & \textbf{Ghi chú} \\
        \midrule
        name.length $\ge 1$ & BV-01 (0) & BV-02 (1) & --- (không max) & Biên dưới bị nhận (BUG-C-01); biên trên khuyết: BV-03 (5.000). \\
        \bottomrule
    \end{tabularx}
\end{table}
\noindent\textbf{Kết luận:} \texttt{name.length} phủ biên rỗng (0), biên hợp lệ (1) và biên trên
khuyết (tên cực dài) --- tất cả quan sát trực tiếp trên bảng danh mục của UI.

\section{AI Gap Analysis}
\begin{table}[H]
    \centering
    \caption{AI Gap Analysis — Feature C (FR-14)}
    \begin{tabularx}{\linewidth}{C{0.5cm} X L{4cm} C{1.8cm}}
        \toprule
        \textbf{\#} & \textbf{Ca/bug bị bỏ sót} & \textbf{Lý do bỏ sót} & \textbf{Bổ sung ở TC} \\
        \midrule
        1 & Tên rỗng/chỉ khoảng trắng được chấp nhận (form không kiểm). & Độ phức tạp (validation không có ở cả UI lẫn backend). & TC-C-DT-02/03, BV-01 \\
        2 & Trùng tên danh mục (không \texttt{UNIQUE}). & Chất lượng prompt (cần đọc kỹ ràng buộc lược đồ). & TC-C-DT-04 \\
        3 & Xóa danh mục làm mồ côi sản phẩm (không khóa ngoại/cascade). & Độ phức tạp nội tại (toàn vẹn tham chiếu, hệ quả ở tab Sản phẩm). & TC-C-DT-08 \\
        4 & Không giới hạn độ dài trên của \texttt{name}. & Độ phức tạp/robustness; BVA thường chỉ nhìn biên đã đặc tả. & TC-C-BV-03 \\
        \bottomrule
    \end{tabularx}
\end{table}
\noindentGhi chú phạm vi: hai khiếm khuyết tầng backend từng ghi nhận ở bản trước ---
thiếu kiểm \texttt{role} khi gọi API, và PUT/DELETE \texttt{id} không tồn tại vẫn trả
\texttt{200} --- \textbf{không quan sát được từ UI} (form login đã chặn non-admin; giao diện
không có nút Sửa và chỉ xóa được hàng có thật). Theo phạm vi functional testing từ UI, hai mục
này nằm ngoài báo cáo.

% ============================================================
%  CHƯƠNG 5: FEATURE D — POOL D (MOBILE)
% ============================================================
\chapter{Feature D — Giỏ hàng (Pool D — Mobile)}

\section{Mô Tả Tính Năng}
Tính năng FR-07 — Giỏ hàng trên app mobile (React Native/Expo) cho phép thêm
sản phẩm vào giỏ, chỉnh số lượng, xóa, rồi thanh toán. Toàn bộ logic nằm trong
\texttt{frontend-mobile/App.js} (giỏ hàng là trạng thái phía client, chưa gọi API cho
tới khi checkout).

\begin{itemize}
    \item \texttt{normalizeQuantity(v)} — \texttt{parseInt}; trả về giá trị nếu hữu hạn và
    $>0$, ngược lại trả \texttt{1}.
    \item \texttt{addToCart(product, qty)} — nếu sản phẩm đã có trong giỏ thì cộng dồn số
    lượng; ngược lại thêm mục mới; luôn hiện \texttt{Alert} ``Đã thêm''.
    \item \textbf{Ô sửa số lượng trong giỏ} — \texttt{onChangeText} đặt
    \texttt{quantity = (parsed>0) ? parsed + 1 : 1}.
    \item \texttt{cartTotal} $= \sum(\texttt{price}\times\texttt{quantity})$.
    \item \texttt{handleConfirmCheckout} — gửi \verb|items = cart.length > 1 ? cart.slice(0,-1) : cart|
    và \texttt{total\_amount = final\_amount ?? cartTotal}.
\end{itemize}

\vspace{4pt}
\noindent\textbf{Trích mã (rút gọn):}
\begin{lstlisting}[caption={Giỏ hàng — App.js (rút gọn)},numbers=none]
// (A) ep so luong hop le
const normalizeQuantity = (value) => {
  const parsed = parseInt(value, 10);
  return Number.isFinite(parsed) && parsed > 0 ? parsed : 1;
};
// (B) o sua so luong trong gio -> luu parsed + 1
newCart[index].quantity = (Number.isFinite(parsed) && parsed > 0) ? parsed + 1 : 1;
// (C) checkout bo mat mon cuoi khi gio > 1 mon
items: cart.length > 1 ? cart.slice(0, -1) : cart,
total_amount: finalAmount,   // = cartTotal (tinh du moi mon)
\end{lstlisting}

\begin{notebox}
Khiếm khuyết trọng tâm: (B) sửa số lượng thành $N$ thì giỏ lưu $N+1$ (off-by-one, thổi
phồng số lượng mỗi lần sửa); (C) khi giỏ có $>1$ món, \texttt{slice(0,-1)} bỏ mất
món cuối khỏi đơn nhưng \texttt{total\_amount} vẫn tính đủ mọi món $\Rightarrow$ khách bị tính
tiền món không được ghi nhận. Ngoài ra (A) âm thầm ép mọi số lượng không hợp lệ
(0/âm/chữ/thập phân) về 1.
\end{notebox}

\begin{notebox}
\textbf{Phạm vi \& quy ước cột kết quả.} Với app mobile, giao diện chính là ứng dụng ---
mọi ca được \textbf{thao tác trực tiếp trên app Expo} (màn hình sản phẩm, giỏ hàng, thanh toán);
cột Kết quả/Trạng thái là hành vi \textbf{quan sát trên màn hình} (số lượng trong
giỏ, dòng ``Thành tiền'', thông báo). Giỏ hàng là trạng thái phía client (chưa gọi API
cho tới khi Xác nhận thanh toán). Chế độ ``sinh viên tự chạy, AI dự đoán từ mã \texttt{App.js}''.
\end{notebox}

\section{Domain Testing}

Áp dụng Domain Testing theo bốn bước (xác định biến $\to$ phân hoạch miền
$\to$ chọn điểm ON/OFF/IN/OUT $\to$ thiết kế ca). Với FR-07, hai biên trọng tâm là
\texttt{quantity} (ranh giới $>0$, phân biệt thêm vào giỏ với ô sửa trong giỏ) và
\texttt{cart.length} (tại 1 --- nơi \texttt{slice(0,-1)} bắt đầu bỏ mất món cuối). Biến đầu ra
\texttt{cartTotal}/\texttt{total\_amount} được đưa vào để lộ mâu thuẫn ``đơn thiếu món nhưng tính
đủ tiền''. Bước~5 tổng hợp ma trận truy vết EC $\to$ TC.

\subsection{Bước 1: Xác Định Các Biến Đầu Vào}
\begin{table}[H]
    \centering
    \caption{Danh sách biến — Feature D (FR-07 Mobile)}
    \begin{tabularx}{\linewidth}{C{0.5cm} L{3.2cm} L{2.6cm} C{1.8cm} X}
        \toprule
        \textbf{\#} & \textbf{Biến} & \textbf{Kiểu} & \textbf{Vào/Ra} & \textbf{Ghi chú} \\
        \midrule
        1 & quantity (trang chi tiết) & string$\to$int & Input & Biên $>0$; \texttt{parseInt} + ép về 1. \\
        2 & quantity (ô sửa trong giỏ) & string$\to$int & Input & Handler lưu \texttt{parsed+1} (lỗi). \\
        3 & cart.length & integer & Trạng thái & Biên tại 1: $>1$ kích hoạt \texttt{slice(0,-1)}. \\
        4 & product.id & integer & Input & Trùng id $\Rightarrow$ cộng dồn số lượng. \\
        5 & cartTotal / total\_amount & integer & \textbf{Output} & Tính đủ mọi món dù đơn thiếu món. \\
        6 & stock / max qty & --- & (thiếu) & Không có mô hình tồn kho/giới hạn. \\
        \bottomrule
    \end{tabularx}
\end{table}

\subsection{Bước 2: Phân Hoạch Miền}
Dấu * = lớp mà hệ thống đúng phải xử lý nhưng mã hiện tại làm sai.

\subsubsection{Biến quantity (thêm vào giỏ)}
\begin{table}[H]
    \centering
    \caption{Phân hoạch miền — quantity (addToCart)}
    \begin{tabularx}{\linewidth}{C{1.4cm} L{6cm} C{1.6cm} X}
        \toprule
        \textbf{EC} & \textbf{Điều kiện} & \textbf{Loại} & \textbf{Ví dụ} \\
        \midrule
        EC-Q1 & Số nguyên dương & Valid & \texttt{3} \\
        EC-Q2 & Mặc định 1 & Valid & \texttt{1} \\
        EC-Q3 & Bằng 0 & Invalid* & \texttt{0} $\to$ ép 1 \\
        EC-Q4 & Âm & Invalid* & \texttt{-2} $\to$ ép 1 \\
        EC-Q5 & Không phải số & Invalid* & \texttt{abc} $\to$ ép 1 \\
        EC-Q6 & Thập phân & Invalid* & \texttt{2.9} $\to$ 2 (cắt) \\
        EC-Q7 & Rất lớn (không có tồn kho/max) & Valid* & \texttt{999999} \\
        \bottomrule
    \end{tabularx}
\end{table}

\subsubsection{Biến quantity (ô sửa trong giỏ) \& cart.length}
\begin{table}[H]
    \centering
    \caption{Phân hoạch miền — sửa số lượng \& số món trong giỏ}
    \begin{tabularx}{\linewidth}{C{1.4cm} L{6.5cm} C{1.6cm} X}
        \toprule
        \textbf{EC} & \textbf{Điều kiện} & \textbf{Loại} & \textbf{Kỳ vọng} \\
        \midrule
        EC-ED1 & Sửa thành $N>0$ & Invalid* & Nên lưu $N$; code lưu $N+1$ \\
        EC-ED2 & Sửa thành 0/không hợp lệ & Invalid* & Reset về 1 (âm thầm) \\
        EC-L1 & Giỏ có 0 món & Boundary & Không cho checkout \\
        EC-L2 & Giỏ có 1 món & Valid & Gửi đủ 1 món \\
        EC-L3 & Giỏ có $>1$ món & Invalid* & Bỏ mất món cuối (slice) \\
        \bottomrule
    \end{tabularx}
\end{table}

\subsection{Bước 3: Chọn Điểm Kiểm Thử (ON/OFF/IN/OUT)}
Hai biên đáng chú ý: \texttt{quantity} tại ranh giới $>0$ và \texttt{cart.length} tại ranh giới
$1$ (nơi \texttt{slice(0,-1)} bắt đầu bỏ món).

\begin{table}[H]
    \centering
    \caption{Điểm ON/OFF/IN — quantity ($>0$) và cart.length}
    \begin{tabularx}{\linewidth}{L{3.4cm} C{2cm} X}
        \toprule
        \textbf{Biên} & \textbf{Giá trị} & \textbf{Kỳ vọng (đặc tả)} \\
        \midrule
        quantity OFF   & 0 & Từ chối/giữ nguyên (code ép 1) \\
        quantity ON    & 1 & Chấp nhận, số lượng = 1 \\
        quantity IN    & 2 & Chấp nhận, số lượng = 2 \\
        cart.length ON & 1 & Gửi đủ 1 món \\
        cart.length +1 & 2 & Gửi đủ 2 món (code bỏ 1) \\
        \bottomrule
    \end{tabularx}
\end{table}

\noindent\textbf{Nhận xét:} biên \texttt{addToCart} xử lý số lượng tương đối ổn (ép về 1),
nhưng ô sửa trong giỏ lại cộng thêm 1 (BUG-D-01); và biên \texttt{cart.length=1$\to$2}
làm lộ lỗi bỏ món cuối (BUG-D-02).

\subsection{Bước 4: Thiết Kế Test Cases — Domain Testing}
\begin{longtable}{C{1.2cm} L{3.2cm} L{3.1cm} L{3cm} C{1.9cm} C{1.4cm}}
    \caption{Test cases — Domain Testing, Feature D} \\
    \toprule
    \textbf{TC\#} & \textbf{Mô tả} & \textbf{Input} & \textbf{Expected Output} & \textbf{Kết quả} & \textbf{Trạng thái} \\
    \midrule
    \endfirsthead
    \multicolumn{6}{c}{\small\itshape (tiếp theo)} \\
    \toprule
    \textbf{TC\#} & \textbf{Mô tả} & \textbf{Input} & \textbf{Expected Output} & \textbf{Kết quả} & \textbf{Trạng thái} \\
    \midrule
    \endhead
    \midrule
    \multicolumn{6}{r}{\small\itshape (còn tiếp)} \\
    \endfoot
    \bottomrule
    \endlastfoot
    TC-D-DT-01 & Thêm với số lượng 3 & qty=\texttt{"3"} & Giỏ có món, qty=3 & qty=3 & \textcolor{green!60!black}{Pass} \\
    TC-D-DT-02 & Thêm số lượng 0 \newline (BUG-D-03) & qty=\texttt{"0"} & Từ chối/giữ 0 & Âm thầm ép qty=1 & \textcolor{orange}{Note} \\
    TC-D-DT-03 & Thêm số lượng âm \newline (BUG-D-03) & qty=\texttt{"-2"} & Từ chối & Ép qty=1 & \textcolor{orange}{Note} \\
    TC-D-DT-04 & Thêm số lượng chữ \newline (BUG-D-03) & qty=\texttt{"abc"} & Từ chối & Ép qty=1 & \textcolor{orange}{Note} \\
    TC-D-DT-05 & Thêm số lượng thập phân \newline (BUG-D-03) & qty=\texttt{"2.9"} & Từ chối/làm tròn rõ ràng & \texttt{parseInt}$\to$2 (cắt) & \textcolor{orange}{Note} \\
    TC-D-DT-06 & Thêm trùng sản phẩm (2 rồi 3) & cùng id & Cộng dồn = 5 & 5 & \textcolor{green!60!black}{Pass} \\
    TC-D-DT-07 & Sửa số lượng trong giỏ thành 2 \newline (BUG-D-01) & ô sửa = \texttt{2} & qty=2 & qty=3 (parsed+1) & \textcolor{red}{Fail} \\
    TC-D-DT-08 & Sửa số lượng trong giỏ thành 5 \newline (BUG-D-01) & ô sửa = \texttt{5} & qty=5 & qty=6 & \textcolor{red}{Fail} \\
    TC-D-DT-09 & Sửa thành giá trị không hợp lệ & ô sửa = \texttt{"x"} & Giữ nguyên/báo lỗi & Reset về 1 & \textcolor{orange}{Note} \\
    TC-D-DT-10 & Xóa một món khỏi giỏ & bấm ``Xóa'' & Món bị xóa & Bị xóa & \textcolor{green!60!black}{Pass} \\
    TC-D-DT-11 & Checkout giỏ 1 món & 1 món & Đơn gồm đúng 1 món & Đúng & \textcolor{green!60!black}{Pass} \\
    TC-D-DT-12 & Checkout giỏ 3 món \newline (BUG-D-02) & 3 món & Đơn gồm 3 món, tổng khớp & Đơn chỉ 2 món, tính đủ 3 & \textcolor{red}{Fail} \\
    TC-D-DT-13 & Thêm số lượng cực lớn \newline (BUG-D-04) & qty=\texttt{"999999"} & Chặn theo tồn kho/max & Chấp nhận & \textcolor{red}{Fail} \\
    TC-D-DT-14 & Checkout giỏ rỗng & 0 món & Chặn (không cho checkout) & Vào màn checkout & \textcolor{orange}{Note} \\
\end{longtable}

\subsection{Bước 5: Ma Trận Truy Vết Độ Phủ (EC $\to$ Test Case)}
Bảng dưới đối chiếu mọi lớp tương đương (Bước~2) với ca kiểm thử phủ nó (kể cả ca BVA), chứng
minh độ phủ đầy đủ cho cả hai nhóm biên \texttt{quantity} và \texttt{cart.length}.

\begin{table}[H]
    \centering
    \caption{Truy vết EC $\to$ TC — Feature D (FR-07 Mobile)}
    \begin{tabularx}{\linewidth}{C{1.8cm} C{3.4cm} X}
        \toprule
        \textbf{EC} & \textbf{Ca phủ (DT/BVA)} & \textbf{Ghi chú} \\
        \midrule
        EC-Q1 & TC-D-DT-01, BV-03 & Số nguyên dương. \\
        EC-Q2 & TC-D-BV-02 & Mặc định 1 (đúng biên). \\
        EC-Q3 & TC-D-DT-02, BV-01 & Bằng 0 $\to$ ép 1. \\
        EC-Q4 & TC-D-DT-03 & Âm $\to$ ép 1. \\
        EC-Q5 & TC-D-DT-04 & Không phải số $\to$ ép 1. \\
        EC-Q6 & TC-D-DT-05 & Thập phân $\to$ cắt phần nguyên. \\
        EC-Q7 & TC-D-DT-13, BV-07 & Rất lớn (không có tồn kho/max). \\
        EC-ED1 & TC-D-DT-07/08, BV-04 & Ô sửa lưu \texttt{parsed+1} (off-by-one). \\
        EC-ED2 & TC-D-DT-09 & Ô sửa giá trị không hợp lệ $\to$ reset 1. \\
        EC-L1 & TC-D-DT-14 & Giỏ 0 món. \\
        EC-L2 & TC-D-DT-11, BV-05 & Giỏ 1 món (gửi đủ). \\
        EC-L3 & TC-D-DT-12, BV-06 & Giỏ $>1$ món (bỏ mất món cuối). \\
        \bottomrule
    \end{tabularx}
\end{table}
\noindent\textbf{Kết luận độ phủ:} 7/7 lớp \texttt{quantity}, 2/2 lớp ô sửa, 3/3 lớp
\texttt{cart.length} đều có ít nhất một ca kiểm thử --- phủ 100\%.

\section{Boundary Value Analysis (BVA)}

BVA cho FR-07 nhắm vào ba biên số lượng / trạng thái: \texttt{quantity} khi thêm vào giỏ
($>0$), \texttt{quantity} ở ô sửa trong giỏ (lưu \texttt{parsed$+1$}), và
\texttt{cart.length} (tại 1). Mỗi biên được kiểm bộ ba \texttt{off/on/on$+1$}; đặc biệt ô sửa lộ
off-by-one (nhập 1 thành 2), còn \texttt{cart.length}$=2$ lộ lỗi bỏ mất món cuối. Bước~3 tổng hợp
ma trận phủ biên.

\subsection{Bước 1: Xác Định Ranh Giới}
\begin{table}[H]
    \centering
    \caption{Bảng giá trị biên — Feature D (FR-07 Mobile)}
    \begin{tabularx}{\linewidth}{L{3.4cm} C{1.8cm} C{1.8cm} C{1.8cm} X}
        \toprule
        \textbf{Biên} & \textbf{Off/dưới} & \textbf{On} & \textbf{On+1} & \textbf{Ghi chú} \\
        \midrule
        quantity $>0$ (thêm) & 0 & 1 & 2 & 0 bị ép về 1; không có biên trên (thiếu tồn kho). \\
        quantity (ô sửa giỏ) & 0/invalid & 1 & 2 & Lưu $N+1$: nhập 1$\to$2, 2$\to$3 (BUG-D-01). \\
        cart.length & 0 & 1 & 2 & $\ge 2$ bỏ mất món cuối (BUG-D-02). \\
        \bottomrule
    \end{tabularx}
\end{table}

\subsection{Bước 2: Thiết Kế Test Cases — BVA}
\begin{longtable}{C{1.2cm} L{3.2cm} L{3.1cm} L{3cm} C{1.9cm} C{1.4cm}}
    \caption{Test cases — BVA, Feature D} \\
    \toprule
    \textbf{TC\#} & \textbf{Mô tả} & \textbf{Input} & \textbf{Expected Output} & \textbf{Kết quả} & \textbf{Trạng thái} \\
    \midrule
    \endfirsthead
    \toprule
    \textbf{TC\#} & \textbf{Mô tả} & \textbf{Input} & \textbf{Expected Output} & \textbf{Kết quả} & \textbf{Trạng thái} \\
    \midrule
    \endhead
    \bottomrule
    \endlastfoot
    TC-D-BV-01 & qty dưới biên (0) \newline (BUG-D-03) & thêm qty=\texttt{"0"} & Từ chối & Ép về 1 & \textcolor{orange}{Note} \\
    TC-D-BV-02 & qty đúng biên (1) & thêm qty=\texttt{"1"} & qty=1 & qty=1 & \textcolor{green!60!black}{Pass} \\
    TC-D-BV-03 & qty trên biên (2) & thêm qty=\texttt{"2"} & qty=2 & qty=2 & \textcolor{green!60!black}{Pass} \\
    TC-D-BV-04 & Ô sửa giỏ nhập 1 \newline (BUG-D-01) & ô sửa=\texttt{1} & qty=1 & qty=2 (parsed+1) & \textcolor{red}{Fail} \\
    TC-D-BV-05 & cart.length=1 & 1 món checkout & Đơn đủ 1 món & Đúng & \textcolor{green!60!black}{Pass} \\
    TC-D-BV-06 & cart.length=2 \newline (BUG-D-02) & 2 món checkout & Đơn đủ 2 món & Đơn chỉ 1 món & \textcolor{red}{Fail} \\
    TC-D-BV-07 & qty biên trên (không max) \newline (BUG-D-04) & qty=\texttt{"999999"} & Chặn/giới hạn & Chấp nhận & \textcolor{red}{Fail} \\
\end{longtable}

\subsection{Bước 3: Ma Trận Phủ Biên (off / on / on$+1$)}
Bảng dưới chứng minh mỗi biên ở Bước~1 đều được kiểm đủ ba giá trị, lấy ca từ cả bảng BVA lẫn
Domain Testing.

\begin{table}[H]
    \centering
    \caption{Phủ biên — Feature D (FR-07 Mobile)}
    \begin{tabularx}{\linewidth}{L{2.8cm} C{1.9cm} C{1.7cm} C{1.6cm} X}
        \toprule
        \textbf{Biên} & \textbf{Off} & \textbf{On} & \textbf{On$+1$} & \textbf{Ngoài / Ghi chú} \\
        \midrule
        quantity $>0$ (thêm) & BV-01 (0) & BV-02 (1) & BV-03 (2) & Không có max $\Rightarrow$ BV-07 (999999). \\
        quantity (ô sửa giỏ) & DT-09 (0/invalid) & BV-04 (1$\to$2) & DT-07 (2$\to$3) & Lộ off-by-one \texttt{parsed$+1$} (BUG-D-01). \\
        cart.length & DT-14 (0) & BV-05 (1) & BV-06 (2) & On$+1$ kích hoạt \texttt{slice(0,-1)} bỏ món cuối (BUG-D-02). \\
        \bottomrule
    \end{tabularx}
\end{table}
\noindent\textbf{Kết luận:} ba biên đều phủ đủ off/on/on$+1$; hai biên ô sửa và
\texttt{cart.length} chính là nơi hai lỗi nghiêm trọng (BUG-D-01, BUG-D-02) biểu hiện.

\section{AI Gap Analysis}
\begin{table}[H]
    \centering
    \caption{AI Gap Analysis — Feature D (FR-07 Mobile)}
    \begin{tabularx}{\linewidth}{C{0.5cm} X L{4cm} C{1.8cm}}
        \toprule
        \textbf{\#} & \textbf{Ca/bug bị bỏ sót} & \textbf{Lý do bỏ sót} & \textbf{Bổ sung ở TC} \\
        \midrule
        1 & Ô sửa số lượng lưu \texttt{parsed+1} — cần đọc đúng phép tính trong handler, không chỉ ``cập nhật số lượng''. & Giới hạn công cụ + độ phức tạp (logic UI mobile). & TC-D-DT-07/08, BV-04 \\
        2 & Checkout bỏ mất món cuối (\texttt{slice(0,-1)}) trong khi tổng tiền tính đủ. & Độ phức tạp (biến đầu ra, chỉ lộ khi giỏ $>1$ món — trạng thái). & TC-D-DT-12, BV-06 \\
        3 & Ép âm thầm số lượng 0/âm/chữ/thập phân về 1. & Chất lượng prompt (dễ coi là ``đã xử lý'' thay vì lỗi UX). & TC-D-DT-02..05 \\
        4 & Không có giới hạn số lượng/tồn kho (biên trên). & Độ phức tạp/robustness; BVA thường chỉ nhìn biên đã đặc tả. & TC-D-DT-13, BV-07 \\
        5 & Lỗi phụ thuộc trạng thái nhiều món (cart.length$>1$) — khó phát hiện nếu chỉ thử 1 món. & Độ phức tạp nội tại (tương tác trạng thái). & TC-D-DT-12 \\
        \bottomrule
    \end{tabularx}
\end{table}

% ============================================================
%  CHƯƠNG 6: AGENT SKILL
% ============================================================
\chapter{Agent Skill}

\section{Mô Tả Skill}
Em xây dựng một Agent Skill tên \texttt{domain-bva-test} (đặt tại
\texttt{.claude/skills/domain-bva-test/SKILL.md}) chạy trong Claude Code. Skill hiện thực hóa
đúng Mục 6 của đề bài: dẫn AI đi qua từng bước của Domain Testing và Boundary
Value Analysis thay vì một prompt chung chung ``sinh test case và tìm bug''.

\subsection{Nguyên tắc cốt lõi}
\begin{itemize}
    \item \textbf{Disciplined AI-First:} skill ép AI thực hiện tuần tự 10 bước, mỗi bước là một
    hạng mục kỹ thuật có thể review được (không nhảy thẳng tới bảng kết quả).
    \item \textbf{Human review:} sau mỗi bước, sinh viên đối chiếu với mã nguồn/đặc tả; skill
    cấm bịa giá trị đầu vào, Expected hay kết quả thực thi.
    \item \textbf{Truy vết nguồn:} mọi Expected phải bắt nguồn từ đặc tả; mọi ``Actual'' phải
    từ thực thi thật trên giao diện (web/app), không đoán mò.
\end{itemize}

\subsection{Mười bước của skill}
\begin{enumerate}
    \item \textbf{Identify variables} --- liệt kê mọi biến vào/ra, kể cả biến ẩn (độ dài, số
    món, ngày, cờ boolean).
    \item \textbf{Characterise domain} --- xác định miền hợp lệ + ràng buộc từng biến.
    \item \textbf{Equivalence partitioning} --- chia lớp valid/invalid, mỗi lớp một mã EC.
    \item \textbf{Identify boundaries} --- đánh dấu biên cho biến có thứ tự/độ dài.
    \item \textbf{Domain matrix} --- phân tích ON/OFF/IN/OUT cho mỗi biên.
    \item \textbf{3-value BVA} --- \texttt{min$-$1, min, min$+$1} và \texttt{max$-$1, max, max$+$1}.
    \item \textbf{Design test cases} --- kết hợp phân hoạch + biên, giả định đơn lỗi + vài ca
    tương tác đa biến.
    \item \textbf{AI gap analysis} --- soi lại ca AI bỏ sót; ghi lý do (prompt/tool/
    độ phức tạp).
    \item \textbf{Emit report} --- xuất báo cáo có mục cho từng bước (thể hiện suy luận, không
    chỉ bảng cuối).
    \item \textbf{Hand off} --- chỉ điền Actual/pass-fail sau khi thực thi; việc báo lỗi và ghi
    AI Audit do người dùng tự gọi các skill \texttt{bug-reporter} / \texttt{ai-audit-logger}
    (các skill không gọi lẫn nhau).
\end{enumerate}

\section{Cách Sử Dụng}
Gọi skill trong Claude Code bằng cú pháp slash-command, truyền feature cần kiểm thử:
\begin{lstlisting}[language=bash, caption={Kích hoạt skill trong Claude Code}]
# Trong phiên Claude Code tại thư mục repo:
/domain-bva-test Áp dụng Domain Testing + BVA cho FR-09 (Mã giảm giá),
đọc backend/server.js và database.js, điền trực tiếp vào main.tex (tiếng Việt),
chế độ "sinh viên tự chạy, AI dự đoán rồi xác nhận bằng thực thi".
\end{lstlisting}

\noindent Metadata của skill (\texttt{SKILL.md}) khai báo điều kiện kích hoạt tự động:
\begin{lstlisting}[caption={Frontmatter của SKILL.md},numbers=none]
---
name: domain-bva-test
description: Apply Domain Testing and Boundary Value Analysis to an
  EShop feature (FR-XX), guiding the AI through every step of the
  technique and producing a reviewable test-case report plus an AI
  gap analysis. Use when the user wants to design test cases for a
  feature, run domain testing / boundary value analysis, or complete
  Section 6 of HW02.
---
\end{lstlisting}

\noindent Skill này đã được dùng để tạo toàn bộ 4 chương Feature A/B/C/D trong báo cáo này
(xem báo cáo AI Audit riêng \texttt{reports/ai-critique-and-audit.md}, Interaction 1--4).

\section{Video Demo}
\begin{itemize}
    \item \textbf{Video 1:} \url{https://youtube.com/...} --- Demo skill \texttt{domain-bva-test}
    end-to-end trên Feature B (FR-09 Mã giảm giá): từ đọc mã nguồn $\to$ phân hoạch miền $\to$
    BVA $\to$ gap analysis $\to$ thực thi API xác nhận bug công thức percent. (Cần quay
    và dán link trước khi nộp.)
\end{itemize}

% ============================================================
%  TÀI LIỆU THAM KHẢO
% ============================================================
\chapter*{Tài Liệu Tham Khảo}
\addcontentsline{toc}{chapter}{Tài Liệu Tham Khảo}

\begin{enumerate}
    \item ISTQB Foundation Level Syllabus (latest edition).
    \item Hardman, P. (2025). \textit{A Post-AI Learning Taxonomy}.
    \item Fuster Rabella, M. (2025). \textit{OECD Education Working Paper No.~338}.
    \item Anthropic (2025). \textit{Building Reliable AI Test Agents} — engineering blog.
    \item DeepEval \& Promptfoo documentation — LLM testing frameworks.
\end{enumerate}

\end{document}